%to have line numbers
%\RequirePackage{lineno}
\documentclass[10pt, letterpaper]{article}      
\usepackage[margin=.1cm,font=small,labelfont=bf]{caption}[2007/03/09]
%\usepackage{endnotes}
%\let\footnote=\endnote
\usepackage{setspace}
\usepackage{longtable}                        
\usepackage{anysize}                          
\usepackage{natbib}                           
%\bibpunct{(}{)}{,}{a}{,}{,}                   
\bibpunct{(}{)}{,}{a}{}{,}                   
\usepackage{amsmath}
\usepackage[% draft,
pdftex]{graphicx} %draft is a way to exclude figures                
\usepackage{epstopdf}
\usepackage{hyperref}                             % For creating hyperlinks in cross references
\hypersetup{pdfborder={0 0 0.4}} %have nice light boxes on refs
 
% \usepackage[margins]{trackchanges}

% \note[editor]{The note}
% \annote[editor]{Text to annotate}{The note}
%    \add[editor]{Text to add}
% \remove[editor]{Text to remove}
% \change[editor]{Text to remove}{Text to add}

%TODO make it more standard before submission: \marginsize{2cm}{2cm}{1cm}{1cm}
\marginsize{1cm}{1cm}{.5cm}{.5cm}%{left}{right}{top}{bottom}   
					          % Helps LaTeX put figures where YOU want
 \renewcommand{\topfraction}{1}	                  % 90% of page top can be a float
 \renewcommand{\bottomfraction}{1}	          % 90% of page bottom can be a float
 \renewcommand{\textfraction}{0.0}	          % only 10% of page must to be text

 \usepackage{float}                               %latex will not complain to include float after float

\usepackage[table]{xcolor}                        %for table shading
\definecolor{gray90}{gray}{0.90}
\definecolor{orange}{RGB}{255,128,0}

\renewcommand\arraystretch{.9}                    %for spacing of arrays like tabular

%-------------------- my commands -----------------------------------------
\newenvironment{ig}[1]{
\begin{center}
 %\includegraphics[height=5.0in]{#1} 
 \includegraphics[height=3.3in]{#1} 
\end{center}}

 \newcommand{\cc}[1]{
\hspace{-.13in}$\bullet$\marginpar{\begin{spacing}{.6}\begin{footnotesize}\color{blue}{#1}\end{footnotesize}\end{spacing}}
\hspace{-.13in} }

%-------------------- END my commands -----------------------------------------



%-------------------- extra options -----------------------------------------

%%%%%%%%%%%%%
% footnotes %
%%%%%%%%%%%%%

%\long\def\symbolfootnote[#1]#2{\begingroup% %these can be used to make footnote  nonnumeric asterick, dagger etc
%\def\thefootnote{\fnsymbol{footnote}}\footnote[#1]{#2}\endgroup}	%see: http://help-csli.stanford.edu/tex/latex-footnotes.shtml

%%%%%%%%%%%
% spacing %
%%%%%%%%%%%

% \abovecaptionskip: space above caption
% \belowcaptionskip: space below caption
%\oddsidemargin 0cm
%\evensidemargin 0cm

%%%%%%%%%
% style %
%%%%%%%%%

%\pagestyle{myheadings}         % Option to put page headers
                               % Needed \documentclass[a4paper,twoside]{article}
%\markboth{{\small\it Politics and Life Satisfaction }}
%{{\small\it Adam Okulicz-Kozaryn} }

%\headsep 1.5cm
% \pagestyle{empty}			% no page numbers
% \parindent  15.mm			% indent paragraph by this much
% \parskip     2.mm			% space between paragraphs
% \mathindent 20.mm			% indent math equations by this much

%%%%%%%%%%%%%%%%%%
% extra packages %
%%%%%%%%%%%%%%%%%%

\usepackage{datetime}


\usepackage[latin1]{inputenc}
\usepackage{tikz}
\usetikzlibrary{shapes,arrows,backgrounds}


%\usepackage{color}					% For creating coloured text and background
%\usepackage{float}
\usepackage{subfig}                                     % for combined figures

\renewcommand{\ss}[1]{{\colorbox{blue}{\bf \color{white}{#1}}}}
\newcommand{\ee}[1]{\endnote{\vspace{-.10in}\begin{spacing}{1.0}{\normalsize #1}\end{spacing}\vspace{.20in}}}
\newcommand{\emd}[1]{\ExecuteMetaData[/tmp/tex]{#1}} % grab numbers  from stata

%TODO before submitting comment this out to get 'regular fornt'
\usepackage{sectsty}
\allsectionsfont{\normalfont\sffamily}
\usepackage{sectsty}
\allsectionsfont{\normalfont\sffamily}
\renewcommand\familydefault{\sfdefault}

%\usepackage[margins]{trackchanges}
\usepackage{rotating}
\usepackage{catchfilebetweentags}

\usepackage{abstract}
\renewcommand{\abstractname}{}    % clear the title
\renewcommand{\absnamepos}{empty} % originally center
%-------------------- END extra options -----------------------------------------
\date{{}\today \hspace{.2in}\xxivtime}
\title{  % remember to have Vistula University!!
  %City Air Makes You Free\\{\large(in some countries, or if you are upper class or financially satisfied)}
  %Cities and Freedom/Agency
  %Freedom and political say across urban-rural\\ 
  %
  %Pereceived
  Personal Agency and Political Efficacy across Urban-Rural 
}
\author{
% Adam Okulicz-Kozaryn\thanks{EMAIL: adam.okulicz.kozaryn@gmail.com
%   \hfill I thank XXX.  All mistakes are mine.} \\
% {\small Rutgers - Camden  % and Vistula University
% }
}

\begin{document}

%%\setpagewiselinenumbers
%\modulolinenumbers[1]
%\linenumbers

\bibliographystyle{/home/aok/papers/root/tex/ecta}
\maketitle
\vspace{-.4in}
\begin{center}

\end{center}

\noindent\textbf{target jou:}\\
  cities if8.2,\\
  ariq if3ish,\\
  some soc j have soc lit: Toennies Simmel Durkheim maybe sth like j or urb soc or eur soc rev?\\

\begin{abstract}
  \noindent We use World Values Survey 7 (2017-2023) across about 50 countries
  % depending on the model like 47-62
  to test urban-rural differences in personal agency and political efficacy.  
%
%   cities are more free. Mostly we find the opposite: in
%   ?? out ?? countries there is mostly no relationship or even in ?? cases
%   positive relationship. Still, in ?? countries cities are indeed more
%   free.
% % none of this robust to country dummies:
% % and we do not find support for it except in Middle East/North Africa. If anything, there is some 
% %   support of the contrary--greater freedom/autonomy in smaller places. But
% %   overall it is a mix and varies byt region and country. 
% %  %
% % East Asia/Pacific stands out with high effect sizes at about .4 or even .5, also
% % controling for other predictors of freedom. South Asia stands out as having large positive effect .3-.4
% % in the smallest category ($<2k$).  
% %  %
% % In particular in many regions or countries there is more freedom v large cities $>500k$ in large towns
% % midsized cities of $100-500k$ and $50-100k$ but slightly less in the smallest
% % places $<2k$.
%
  We find that there isn't more personal agency in urban v rural as
  hypothetised--there is mostly no difference. But political
  efficacy is significantly higher in urban v rural in rich countries as hypothetised in ``left
  behind ruralities'' hypothesis. And as Toennies hypothetised 
  those doing well financially are ``active and alive'' in cities or as we find here have
  greater personal agency (but not political efficacy). Effect sizes are
  considerable at about ??. As in any correlational research results may not be
  causal. 
    
%  Within-country urban-rural freedom/autonomy differences are mostly
%  insignificant, but some effect sizes tend to
%  be substantial, about 0.3-0.6 % and as large as around 1
%  on 1-10 freedom/autonomy scale. %fist results table by country
%  %
%  %
%   %Toennies argued that ``city life and Gesellschaft down the common people to decay and death.''%  We have not measured
%   % decay and death here, but
%   %Likewise, we find that
%   And lower class, financially dissatisfied, and the poor 
%  are more free/autonomous outside of the cities.   

%  % % i guess not, bc without country dummies: 
% % The whole sample pooled effect sizes of financial satisfaction and class interactions with
% %urbanicity are substantial of about 2 steps on 1-10 income scale.
% %%.12 to .16 %interactions in app table underlying body marginsplot on finSat and class similar to .13 on 1st inc dummy v base so 2 steps
\end{abstract}
\vspace{.15in} 
\noindent{\sc World Values Survey, freedom, agency, autonomy, political efficacy, urbanism, urban, rural,
  urban-rural, cities%XXX TODO add to ebib as keyword PAPER-CODE-NAME and tag with ebib keywords 
}
\vspace{.25in} 

\begin{spacing}{1.4} %TODO MAYBE before submission can make it like 2.0
\rowcolors{1}{white}{gray90}

\addtocontents{toc}{\protect\setcounter{tocdepth}{-2}}


\noindent City life and Gesellschaft down the common people to decay and death. Ferdinand Toennies\\
%\citep[p. 231]{tonnies57}

\noindent In the city and therefore, where general conditions characteristic of the
Gesellschaft prevail, only the upper strata, the rich and the cultured, are
really active and alive. Ferdinand Toennies\\ %\citep[p. 227]{tonnies57}.


City air makes you free%(City air brings freedom) ``City air makes men free
\footnote{(Stadt Luft macht frei) \citep[p. 12]{park84}}{--the saying originated in the Middle Ages, and it meant
  freedom from feudalism in cities, non-feudal islands in a sea of feudalism
  \citep{harvey12}.} 
%The parallel is somewhat a stretch--and the title is somewhat a hook bait.
%The original saying is from the Middle Ages and meant freedom from feudalism in cities.
Feudalism has long been gone, but we may actually be living under feudalism
again, so called techno-feudalism
\citep{varoufakis24,morozov2022critique}. While %there is no reason to believe that
modern cities aren't clearly more free from techno-feudalism,\footnote{Cloud
  fiefdoms, rent, and digital serfdom through platforms like Google and Amazon
  ``enslaves'' urbanites and ruralities, but still perhaps city may be more free
in a sense that there is more choice and opportunity say with local
traditional on the ground commerce.} as Middle Ages cities used to
be more free from feudalism, cities are still often seen in the literature and
popular culture as more free than ruralities.
%leonie didnt like: 
%especially for non-conformist types such as LGBTQ.
%
%
%
 Hence, it'd be already a finding worth reporting that people don't have more
 freedom/agency in cities, because often the expectation is that they do.
 %there are theories, reasons, and arguments that they do. 

%guess could add on escape from freedom fromm
Freedom/agency is an end-in-itself, a desirable outcome in just about any
respect, except perhaps responsibility. In the Middle Ages there was little freedom, but
also little responsibility--for instance the poor and the slaves were not
responsible for their lot--their situation was God given. Now the poor have to worry about
being poor and about it being their fault--everyone is free after all and their
achievements and experiences are of their own making.\footnote{For instance,
  people in the US feel relatively free and so they blame the poor for their
  lot, while Poles feel less free and also blame the poor less. %in addition to
                                %free var could add poor are lazy etc from wvs  
%  Still, blaming the poor
%for their lot is often more prevelent in less free countries such as the US than in
%more free countries such as Mexico or Colombia.
}

%again not robust to country dummies
 % We test hypothesis that cities are more free, and we do not find
 %  support for it except in Middle East/North Africa. If anything, there is some 
 %  support of the contrary--greater freedom/autonomy in smaller places.

 We measure freedom as (perceived) personal agency, which we discuss first, and
 then we turn to our other main variable of interest, (perceived) political
 efficacy.  

\section{Preceived personal agency %Freedom/autonomy defined
}
 
% Some authors use term ``freedom'' \citep{aok_free_from_to} and some ``autonomy''
% \citep{steckermeier22}--we use the terms interchangeably, but prefer
% ``freedom'' as it seems to sound and fit better in this context.
 
What is freedom/agency?\footnote{Here we focus on personal as opposed to
  societal freedom and follow \citet{aok_free_from_to} in this paragraph.} To be free is to be able to choose, to control and
direct one's own life. Freedom is having control--a person is free is she is her
own master. 
 It could be labeled ``psychological freedom:'' ``Freedom means a degree of harmony
between basic motives and overt behavior'' \citep[Bay 1970, p.83 cited in][]{aok_free_from_to}.
 It is ``capability to choose,'' ``awareness of alternatives'' and ``courage to
 choose'' \citep{veenhoven00_b}.\footnote{The degree of perceived control over choice is called ``locus of control'' in
 psychology and \citet{verme09} argues that WVS item used here measures locus of
 control in addition to freedom.}
 Clearly, there are more choices in city v
ruralites, so at least in a sense of choice there is more freedom. But then too much choice may
 also backfire and even lead to a paralysis \citep{schwartz04}.

Freedom/agency is a broad concept and could be sub-divided into dimensions--\citet{hitlin09} reviews the concept in sociology. 
See also a recent scale \citep{agencyScale25}. Another useful review is
\citet{kouba16} and also relevant here as it points to our WVS measure.


There are theoretical reasons to expect more freedom in cities v ruralities. 
In rurality there is mechanical solidarity--rural folks function like parts of a
simple machine--groups are held together through homogeneity and similarity:
similar lifestyles, roles, shared values and morality \citep{durkheim97}. Hence,
 ruralities force social
conformity not conducive for individual freedom. On the other hand, organic
solidarity in cities allows for more differentiation and individualism. %and individuation. %(Jung)
 Urbanites are also more tolerant \citep{tuch87}.
 % 
% 
Simmel observed that as urbanism overloads a person it forces to
distance oneself from the environment \citep{simmel03}, which also may be
conducive for differentiation, individualism,  %and individuation.
 and then perhaps freedom. On the other hand, urban hunkering down, withdrawal,
 and blase attitude \citep{simmel03} are not agentic.  
%
% Likewise, there would be more Gemeinschaft in ruralities--more collective and
% social that may restrict personal freedom; and more Gesellschaft in cities that
% may promote individualism (but also hostility and aggression). 
% 
Cities provide freedom from traditional rural social constraints--Gesellschaft v
Gemeinschaft--deviation from community norms is more likely to be reprimanded
and ostracized in rural v urban.
An urbanite may be more internally driven rather than
externally constrained rural folk \citep{yamagishi12}. %rubia wa saying sth like that in one of our papers too: 
 City provides more opportunities for alternative jobs and social relations, so
 urbanite is more free by being less concerned about social exclusion as she has
 plenty of alternatives \citep{yamagishi12}. Also see 
 discussion in \citet{luca2023progressive}. %just srch for 'free' or 'freedom' 
 
 One mechanism is higher social pressure for conformity in ruralities--cities
 provide anonymity--indeed concern for moral reputation is greater in rural than urban areas
 \citep[Danielson et al 2023 cited in][]{aokOttCh26}. The higher the population
 density, the lower the frequency of social comparison measuring social
 comparison with searches on Google \citep[][]{aokOttCh26}.   
 
 But there is also theory suggesting less freedom in cities.
 %
 City %(and the US)
 is stressed, hectic, anxious, materialistic, consumerist
 \citep{aokCityBook15}, and overstimulating \citep{simmel03,lederbogen11}--in such conditions freedom/autonomy is
difficult.
 %
  Steve Pile in his colorful writing  argues that old cities can be a weight on its inhabitants with their
 history and past \citep{pile05,pile05B,pile99}, and  old cities carry
 melancholia \citep{pile05B}. Weight of history and melancholia can inhibit how
 free one feels.

  
Many studies approach freedom as a predictor of subjective wellbeing, but freedom
is also an end-in-itself \citep{aok_free_from_to}, as recognized by 
\citet{fromm13,fromm12,fromm92,fromm64,fromm44,fromm94} and \citet{sen00}.  The Universal
 Declaration of Human Rights states that freedom is a human right: 
 freedom for the realization of one's human potential (art. 26, 29). 

There are studies on freedom using WVS variable that we use in our study  
\citep[e.g.,][]{aok_free_from_to,pitlik2016free,nikolaev2016give}, but they do not examine
urbanicity. 
There are many studies on the effect of urban-rural on subjective wellbeing (SWB)
\citep[e.g.,][]{sorensen24,sorensen14,knies26,aok21}, but the effect of
urban-rural on freedom has not been studied much. 


%leonie
\citet{lolle2023differences} finds similar non-differences in agency across
different sized communities in Denmark. 
\citet{atherton2024rural} find that psychological well-being--which entails
 agency as a subcomponent--is lower in the most rural places compared to the
most urban places in the USA; however, these effects turn insignificant once
socio-demographic factors are controlled for. 
\citet{kusnali2025rural}  find that women in Indonesia have more decision making
 agency in many (household expenditures, food expenditures, time allocation to
social activities) but not all (contraceptive use) aspects of daily life in
urban compared to rural areas.
 Baage et al find no urban-rural difference in peoples
perceived reproductive agency, i.e. how much freedom of choice and control they
have over family planning, in Ethiopia, Kenya, and Zimbabwe
\citep{baage2024attitudes}  and sub-Saharan Africa \citep{baage2026reproductive}.

 
\section{Perceived political efficacy}

We turn to political efficacy. %AI:
It's a person's belief that she has a say in politics, a perception that
goverment will respond to one's voice. 
 
%\citet{lord2023autonomy} is a summary of their themed issue on ``Autonomy and
%independence in urban environments.''  Rather suburban than rural contrasts for
%urban, and more on culture and life cycle than space.
 %
As the world continues to urbanize, ruralities are disappearing 
%, and so there is also  ``the perceived disappearance of rural childhood and the values
%associated with it'' in Czech republic
 \citep{pospvech2024loss}. Rurality in
 the developed world is increasingly a periphery, loosing power and voice v
 urban centrality, being abndoned and disrespected--there is rural disconent, resentment, and shame. See also
 \citet{hansonCityJournalautumn15,fullerNYT17monD}. This would suggest that
 ruralities in rich countries may have less autonomy--but that's rather
 because ruralities have been abandoned and disrespected, not because they are
 inherently less free or agentic. 
 %!!!opposite to our iniitial country level hypothesis that "This difference is more pronounced in less developed countries"
 %
 %
 %what makes one a real rural person: in a classical study, Bell (1994) identified
 %localism (a history of life in rural areas), ruralism (an experience with a rural kind of job), countryism (knowledge of country life) and communalism (interest in things that express a ties to the
 %community) as key elements of belonging to a place.
 %Peripherality, in other words, is defined through long distances and poor infrastructures but
 %does not have to be experienced through them.
 %
 %
 %And ruralities became digitized--children do not play outside in natura that
 %much anymore, but instead on screens just like their urban counterparts. And
 %autonomy is lost to screens--kids are dependent or even addicted to them.
 %
 %
 In general, while there has been a loss of power and respect in rural v urban, not
 necessarily the autonomy has been lost, ruralities still have a good deal of agency. 

 Like in Czech Republic, in the US there is rural resentment, too, e.g., that
 policy is done to them, not with or for them \citep{diamond23}.  And in
 addition there is racial resentment where rural Whites resent urban minorities
 and immigrants getting better treatment than them \citep{diamond23}. %p505
 % handouts and skip the line
 %
 %ruralities more place attachment and nature conservation
 And while there is a reliance on comunity and tight connection, there is also
 self-reliance, and desire for autonomy and freedom as a core value
 \citep{diamond23}.


%leonie 
Rural areas have often been characterized as 'left-behind' \citep{ulrich2021people} or as 'places that
don't matter' \citep{rodriguez2018revenge}, %cited like 4k!!
 reflecting the notion that they are second-tier areas in which opportunities 
 are limited or deteriorating. 
  %
When people's opportunities are constrained, their freedom to exercise agency is
likewise restricted. Exercising agency requires access to real and substantively
different opportunities from which individuals can choose in pursuing lives that
are consistent with their own goals and values. Limited opportunities may
therefore undermine both, people's perceived ability to shape their own lives in
accordance with these goals and values and their perceived ability to shape the
conditions in which their lives are embedded, such as institutions,
infrastructure, and the built and natural environment.
 %
We therefore expect people living in 'left-behind' rural areas to perceive
themselves as having less influence over the conditions in which they are
embedded and that shape their lives, and consequently to perceive lower levels
of political and personal agency.
 


%leonie
Only few studies have investigated urban-rural differences of political
efficacy, suggesting people from rural areas and small towns show lower levels
of political efficacy compared to urbanites, \citep[e.g.,][]{rowland2024political} for
the US and \citet[e.g.,][]{luukkonen2022urbanisation} for Europe.
 There is more political efficacy in urban areas and least in
rural in Europe \citep{garcia2024they,luukkonen2022urbanisation}, and USA
 \citep{rowland2024political}.  
 %
 And no one has investigated whether this relationship is universal across
 different contexts (such as high/low GDP).
 
 %leonie
 \citet{margalit2025cultural} %looks like may be important study
 examine political agency in relation to rurality,
conceptualizing ``rural resentment'' as the perception among rural citizens that
the political system gives people like them little or no say in what the
government does. Across 10 European countries, they consistently find that rural
resentment is associated with a higher likelihood of voting for populist
parties, while in the United States it is associated with a higher likelihood of
having voted for Trump.
%
Political efficacy has been shown to be positively linked to political
engagement \citep{accetta23}.\footnote{Self-efficacy and political efficacy have
  been shown to have opposing effects on the justification of terrorism and
  political violence: people who perceive themselves as highly self-efficacious
  and those who feel they have little political efficacy are less likely to
  justify such actions \citep{julkif2022self}.}
 

%glue somewhere
Indian females have more autonomy in city v rural, both in personal and 
 household decisions \citep{mishra2014relative}.
 





\section{Hypotheses}

These are in terms of both personal agency and political efficacy:

$H_1:$ People in cities are freer than people in rural areas.\\

$H_2:$ This difference is more pronounced in less developed countries\\
%i guess bc diffs are larger as per my addis abeba disst argument\\
%but yes with perc urban--more urban more c, more freedom in city: in fig \ref{t-urb}

$H_3:$ The "more" of freedom in cities mostly stems from people satisfied with
finances (following the Toennies quote)\\

%Add on political efficacy: abondoned rural in the rich world theory is there


\section{Data and model}

 World Values Survey (WVS) is a fantastic and free dataset available from 
 \url{https://www.worldvaluessurvey.org}. It is a great source of data as
 compared to $\$30k+$ closed source Gallup data. It is helpful to use open
 source public data amidst corporatization of research and academia
 \citep{mills2012corporatization,cox2013corporatization,millsNYT12fa,CatropaNYT20feb8,schmidlinNYT15oct10}. 
 %
WVS is representative of each country. Countries in the sample and sample sizes are shown in appendix.



Urban-rural is a gradient, not a dichotomy \citep{aok-ls_fisher16}, but here for
simplicity we use 3 or 4 collapsed categories and focus on the endpoints of the spectrum: $>500k$ v small $<10k$ or $<20k$ (separately for each
country). $>500k$ represents city--indeed it could even be bigger than 500k--urban
unhappiness starts around several hundred thousand \citep{aok-ls_fisher16}.

In appendix we do robustness checks droping countries with fewer than 50
observations per category.
%We dropped countries where there were fewer than 50 observations in cities $>100k$
%And in appendix as a robustness/sensitivity check with fewer than 50 observations in cities $>500k$


\subsection{Personal Agency}

Autonomy/freedom variable reads:
``Some people feel they have completely free
choice and control over their lives, while other people feel that what they do
has no real effect on what happens to them.  Please use this scale where 1 means
'none at all' and 10 means 'a great deal' to indicate how much freedom of choice
and control you feel you have over the way your life turns out.''

% Validity and reliability of the measure seem fine and the variable has been used successfully
% in the literature \citep[e.g.,][]{verme09,inglehart08,veenhoven00_b,aok_free_from_to} but still
% the question may be interpreted slightly differently in different
% languages/cultures, and some caution would be useful. This is also one reason for
% within-country analyses that we start with in the next section. 

All variables including controls are listed in appendix and there are
also shown distributions of all the variables. 
We are controlling for %age, age$^2$, male,
married, employment type dummies (full/part time, self-employed, retired,
housewife, student, unemployed, and other),
financial satisfaction,\footnote{Financial satisfaction is the key--not everyone
  needs the same amount of money to be financially satisfied and feel free.
  It is measured as ``How satisfied are you with the financial situation of your household?''} health, kids,
and importance of religion.

Per controls we follow \citet{pitlik2016free,nikolaev2016give} except that we use financial
satisfaction in most specifications instead of income as it arguably predicts
 freedom more than income and indeed in our sample correlates higher, and also
 we don't use their country level democracy index, and use PCGDP in appendix
 only. Eventually we decided against age and and age$^2$, as it is
 rather urbanicity predicting age composition than the other way round, hence we
 would risk overcontrol bias. Neither we include gender as it likely differs
 less than age across urban-rural, and again rather predicted by than predicting
 urban-rural.  

We were tempted to add trust, crime victimization and fear of crime, as these
correlate with urbanicity and may predict freedom. 
%and ineed inclusion of these yields larger coefficients.
 %
 % ``Have you been the victim of a crime during the past year?'' and
 % ``In the last 12 months, how often have you or your family: Felt unsafe from
 % crime in your own home''  
 %
% Yet we only include trust and crime victimization and fear of crime as last
% oversaturated model,
But eventually we do not include them because those are rather caused by urbanicity, not the
other way round, and hence may lead to overcontrol bias:
\begin{quote}
What we really need, to identify/estimate X$->$Y, is to control for other
determinants of Y that are also antecedents of X (so, W$->$X). If we include
controls where the relationship goes the other way (X$->$W), we will exacerbate
bias in our estimate of X$->$Y. (The relevant term for that situation is
'overcontrol bias') \citep[][p. 5]{bartram24}.
\end{quote}


\subsection{Political Efficacy}

 %leonie
 The  political efficacy item reads ``how much would you
 say the political system in your country allows people like you to have a say
 in what the government does?''
 % This was a 5-point scale with a higher number
 % signaling less efficacy. This question was fielded from 2018 onward, so
 % countries whose fieldwork was completed before the question was added (except
 % Argentina where the question was fielded in 2017) were dropped from this
 % analysis due to missing data.
 There are quite a few countries where the question was not asked:
 the United States, Germany, South Korea, Greece,
 Turkey, Thailand, China, Chile, Ecuador, Serbia, Romania, Russia, and
 Tunisia.\footnote{Although
 Tunisian respondents were asked this question, the data for Tunisia were not
 curated properly and so could not be analyzed in this
 study. \url{https://onlinelibrary.wiley.com/doi/full/10.1111/ssqu.13120}} 
 
% leonie
Controls. Education: more educated people more efficacy
\citep{emmenegger2015labour,garcia2024they,marx2016unemployed,rowland2024political}.
Employment: compared to paid employment people in education \textbf{???}  more political efficacy,
permanently sick/disabled less \citep{luukkonen2022urbanisation}, unemployment or
precarious employment are negatively linked to political efficacy
\citep{emmenegger2015labour,marx2016unemployed}. Financial situation: People who
do well (making ends meet), have more political efficacy
\citep{luukkonen2022urbanisation,marx2016unemployed}. Higher income $=$ more
political efficacy \citep{emmenegger2015labour,rowland2024political}. 
 Health: positive for political efficacy
 \citep{garcia2024they,gidengil2024healthy,shore2019health}.
 Religiosity: more religious people have less political efficacy \citep{marx2016unemployed,rowland2024political}.
 %
 Why these controls: There are indications for urban-rural divides
 in education \citep{van2021urban,zahl2024spatial},
 employment  \citep{e26,ananian2024employment},
 financial situation
 \citep{kpatcha2026unlocking,meloni2024rural,finnemann2024urban},
 health \citep{mackinnon2023mapping,menashe2022shifting},
 and religiosity 
 \citep{balazka2020mapping,meulemann2020secularization}.



\section{Results}

Means of personal agency and political efficacy across urban-rural and level of
development are in table \ref{m0}. It is worth highlighting that a finding of even no difference
 in personal agency beteen urban and rural would be worth reporting as we have hypothetised a
positive relationship.  

Overall personal agency is the lowest in the largest places,
including in less developed countries, and 
the difference being especially large for high development countries.
But  in upper middle it is actually highest in cities. But overall, the
differences are not that large and go away when controlling for predictors of freedom.

Political efficacy is also overall slightly lower in largest cities. But as opposed
to freedom differences are large by level of development. In less developed
countries cities have much less political say; in upper middle less, but in high
almost the opposite of low/lower, and that last difference holds when controling
for predictors of political say. Margins are plotted in figure \ref{m-fps}.

While the confidence intervals are wide,
the interaction between $>$500k and high development countries is significant v
the base case of $<$20k and low/lower development countries. The results for
personal agency are insignificant. 


\begin{table}[H]\centering\footnotesize
\caption{\label{m0} means}
\begin{tabular} {@{} lrrrr|rrrrr @{}}   \hline 
       & \multicolumn{4}{c}{free}          &\multicolumn{4}{c}{political say} \\
place  &overall&low/lower middle&upper middle&high&overall&low/lower middle&upper middle&high\\ \hline
$<$20k &7.20   &7.11     &7.26        &7.24&2.96   &3.03     &3.06            &2.67  \\
20-500k&7.27   &7.21     &7.21        &7.36&2.88   &2.95     &2.83            &2.87  \\
$>$500k&7.12   &6.96     &7.40        &6.98&2.89   &2.67     &2.86            &2.99   \\ \hline
\end{tabular}\end{table}

\begin{figure}[H]
 \includegraphics[height=6.5in]{m-fps.pdf}\centering
\caption{\label{m-fps} }
\end{figure}


\subsection{Interactions with financial satisfaction}

% Again, the initial results with pooled sample for all the countries  (in appendix) turned out to be all insignificant when including country dummies, hence we started with results within each country in the previous
% table. 

% But there is one significant result for the pooled sample on urbanicity when
% controlling for countries--when urbanicity is interacted with financial
% satisfaction, class, and income. While earlier we just used OLS within each country,
% here we use multilevel modeling (MLM) with random intercept at country level.

 Results are visualized in figure \ref{m-body}.\footnote{We focus here on visual
  interpretations using marginsplot, but underlying regressions are in appendix. %1.5 Main results from paper: Interactions with financial satisfaction, class, and income
The relevant significant interactions are for smallest places $<10k$ (against
the base case of $>500k$ with top category(ies) for financial satisfaction in
first column a for middle financial satisfaction (v bottom base case) at -0.12
and the top category at -0.13. % Likewise, in column b for class interactions,
% there is significant interaction of $<10k * Upp mid/upper$ at -0.16. Thus, folks
% of both  high financial satisfaction and class are less free in smallest areas v
% the largest areas. In last column c with income, there are no significant
% interactions except unexpected top income quintile, but other insignificant
% patterns are in line with financial satisfaction and class: the two bottom
% deciles are less free in the largest places.
 For political efficacy in panel b results are insignificant. 
}


In panel a by financial satisfaction those in
top and also middle 3 categories increase from $<20k$ to $>500k$.
% 
% 
In other words, freedom differences are magnified in largest cities for
financial satisfaction--the financially satisfied enjoy higher freedom premium
over financially dissatisfied  in cities. The results for political efficacy are
 insignificant. 



\begin{figure}[H]
 \includegraphics[height=6.5in]{zz.pdf}\centering
\caption{\label{zz} }
\end{figure}



% We also separate these results by country, however we leave
% it for the future research as sample sizes run small for interactions.

We also separate results 
 by region and income group for financial satisfaction and  postpone these results to appendix\footnote{Here just
real quick: Financial satisfaction panel a results are stronger in East
Asia/Pacific and Sub-Saharan Africa (due to sharp increase for the top tertile), and show similar pattern in Europe/Central Asia except that the
middle tertile for financial satisfaction is stronger than the top. Latin
America/Caribbean and North America shows the opposite pattern where  the bottom tertile has the
steepest slope, and in Middle East/North Africa middle and top tertiles even reverse
to negative. in South Asia the slopes are parallel--no differences there. By
income group: High income is similar to the overall pattern.}

\section{Discussion and Conclusion}

% In the pooled sample without interactions there are no significant urban rural
% differences when controlling for country dummies. But there are some significant
% results separately
% by country. The hypothesis that cities are more free does not find support
% except in 16 countries out of 64.

It is worth highlighting that finding no urban-rural effect for personal agency
in most countries is a finding worth reporting--the expectation from theory and
literature is that the bigger places have more freedom. %And in 6 countries we find a positive effect. 

But political
  efficacy is significantly higher in urban v rural in rich countries as hypothetised in ``left
  behind ruralities'' hypothesis.

We find support for Toennies' idea:
 ``in the city and therefore, where general conditions characteristic of the
Gesellschaft prevail, only the upper strata, the rich and the cultured, are
really active and alive'' \citep[p. 227]{tonnies57}.
 There are significant interactions in the pooled sample for financial
 satisfaction. 

  
Like in \citet{aok24swls} urbanites probably actually have in a way fuller and
more free life, but perhaps
they realize more that they miss out relatively to other urbanites, and so they
feel relatively less free. One mechanism could be that people tend to compare
upwards \citep{frey02s} and there are more people to compare to in urban areas,
inequality ladders are taller and the pond is bigger--small fish in big pond
(city) seem to be less free than big fish in a small pond (rural area).

In poorer/less economically
developed countries there are more differences in freedom across rural-urban, mostly less freedom in smaller
places\footnote{(but sometimes more)}--this makes sense--smaller places in
developing countries can be very traditional and very poor--freedom is difficult
in such places. In richer countries, on the other hand, one can already be free
just about anywhere, no need for a large city. This is a similar finding as in 
literature with respect to happiness \citep{aokCityBook15}--only in
the very poorest places of all such as Sub-Saharan Africa, happiness is greater in
larger cities as life is unbearable and lacking necessities such as shelter and
clean water in ruralities. Everywhere else urbanites tend to be less happy \citep{aok21}.


%  instead \ExecuteMetaData[../out/tex]{ginipov} do \emd{ginipov}

% \begin{figure}[H]
%  \includegraphics[height=3in]{../out/gov_res_trust.pdf}\centering\label{gov_res_trust}
% \caption{woo}
% \end{figure}


%TODO !!!! have input here aok_var_des



% %table centered on decimal points:)
% \begin{table}[H]\centering\footnotesize
% \caption{\label{freq_im_god} importance of God}
% \begin{tabular} {@{} lrrrr @{}}   \hline 
% Item& Number & Per cent   \\ \hline
% 1(not at all)&    9,285&  9\\
% 2&    3,555&        3\\
% 3&    3,937&        4\\
% 4&    2,888&        3\\
% 5&    7,519&        7\\
% 6&    5,175&        5\\
% 7&    6,050&        6\\
% 8&    8,067&        8\\
% 9&    8,463&        8\\
% 10&   52,385&       49\\
% Total&  107,324&      100\\ \hline
% \end{tabular}\end{table}


% % Define block styles
% \tikzstyle{block} = [rectangle, draw, fill=black!20, 
%     text width=10em, text centered, rounded corners, minimum height=4em]
% \tikzstyle{b} = [rectangle, draw,  
%     text width=6em, text centered, rounded corners, minimum height=4em]
% \tikzstyle{line} = [draw, -latex']
% \tikzstyle{cloud} = [draw, ellipse,fill=black!20, node distance = 5cm,
%     minimum height=2em]
    
% \begin{tikzpicture}[node distance = 2cm, auto]
%     % Place nodes
%     \node [block] (lib) {liberalism, egalitarianism, welfare};
%     \node [block, below of=lib] (con) {conservatism, competition, individualism};
%     \node [cloud, right of=con] (ls) {well-being};
%     \node [block, below of=ls] (cul) {genes, culture};
%     \node [b, left of =lib, node distance = 4cm] (c) {country-level};
%     \node [b, left of =con,  node distance = 4cm] (c) {person-level};
%     % Draw edges
%     \path [line] (lib) -- (ls);
%     \path [line] (con) -- (ls);
%     \path [line,dashed] (cul) -- (ls);
% \end{tikzpicture}


%PUT THIS NOTE, polish and put to /root/author_what_data --ALWAYS
%stick here stuff as i run it!!! maybe comment out later...


%TODO: have separate som-r.tex as opposed to having it below; and in paper say
%see online appendix!

\setcounter{section}{0}
 \part*{\Huge ONLINE APPENDIX (SOM)}
 \textbf{[note: this section will NOT be a part of the final version of
   the manuscript, but will be available online instead as SOM 'Supplementary Online
   Material']} %hence everything below

\tableofcontents
\addtocontents{toc}{\protect\setcounter{tocdepth}{3}} 





\section{More recent results}

\subsection{Literature}

\subsubsection{Freedom/personal agency}
{While \citet{ester93} is not about urban-rural freedom specifically,  it is
 indirectly related.  
 \citet{ester93} uses similar dataset to ours, EVS/WVS, and tests yuppie value syndrome (yuppies are young urban
 professionals), and finds that yuppies tend to hold liberal social and ethical
 views, but are not: %as hypothetised:
 economically conservative, absent from
 politics, and don't have distinctive lifestyle emphasis. 
}

\citet{yamagishi12} claims to test
 our ``city free air'' hypothesis, % : the social constraints prevalent in rural life are
 % weaker in metropolitan areas, freeing metropolitan residents from pressure to
 % suppress their pursuit of individual goals.
 but \citet{yamagishi12} test is very different
 from ours--they only show that urbanites are more likely to chose extravagant pens
 than rural folks. Still, the study is useful in at least 2 ways. First they
 show that indeed there is more freedom in city (using extravagant pens test),
 and they provide theory and literature summarized in the next paragraph.   
 

Perception of freedom can be counterintuitive and deceptive--a place may seem free, or indeed
actively advertise itself as an oasis of freedom, but actually not be that free
after all. 
 %
 %City, just as the US appears the land of the free, but it is not.
 %Freedom/autonomy can be counterintuitive--for instance
 The US advertises itself
as ``the land of the free'' and yet people in the US feel less free at 7.7 than
Mexicans and Colombians at above 8.1 (WVS Wave 7). Likewise, city may seem more free than
ruralites, but as this paper finds, usually that's not the case.
  %
%But we also have the case why
% There would be less freedom/autonomy in the city, just as
% there would be less in the West (especially the US) v Latin America. Mexico,
% Colombia and Venezuela are some of the most free countries in the World
% \citep{aok_free_from_to,colQolSwb25}. 
 %
 %
In a city v rural (or in the US v Latin America) there is often  
 more social pressure \citep{colQolSwb25,unhappyBog26}\footnote{Also see a fantastic book by \citet{frank12} on
   keeping up with the Joneses in the US.} and so an urbanite is more of
 an ``automaton'' (to use Fromm's term)
% Also, an American may be confused or even paralyzed because there is toomuch variety (as Schwartz (2004) suggested) and accordingly she does not feelfreedom to.  Yet another reason may be that a Mexican does not knowthat she is not free relative to an American–she feels more free relative to other Mexicans than an American feels free relative to other Americans.
 \citep{aok_free_from_to}.%footnote 7 



Money is important for personal freedom/agency. Indeed as many financial gurus on YouTube would
tell you freedom is the only or main reason to accumulate wealth. For instance, the FIRE movement (Financial Independence, Retire Early) is
all about freedom/autonomy. Hence in our analyses we will focus on income, class
and financial satisfaction.


\subsection{Original (and before the final) results from the paper postponed here}
 
A natural first approach would be to find overall pattern on pooled complete
dataset across all countries to see if there are universal patterns or at least by income group or
region. Indeed this is how we have started, but we postpone these results to the
appendix as we have learned through our sensitivity/robustness checks that
inclusion of country dummies yields all urban-rural results insignificant.
Hence we simply start with the results by country.

Financial satisfaction is the strongest correlate of freedom (0.33), 
stronger  than health (0.21), class (0.14), and income (0.09).
 Financial satisfaction is also related to class (0.26) and income (0.34). 



{\bf By country}

Results are in figure \ref{fig-regCC1mod}. 
Is there less freedom in smaller areas as often argued? Are cities more free?
No. Most countries show insignificant urban-rural differences at  $<0.05$ level
of significance.  Again, it is already a finding worth reporting that people
don't have more freedom/autonomy in cities, because there
are theories, reasons, and arguments that they do.


\begin{figure}[H]
 \includegraphics[height=5in]{fig-regCC1mod2}\centering
\caption{\label{fig-regCC1mod2} A more final version, where country count goes
  down from 66 to 51
   when dropped if gt500k missing; same pattern as
  earlier: most insignificant, some positive but most negative.} 
\end{figure}


\begin{figure}[H]
 \includegraphics[height=5in]{fig-regCC1mod}\centering
\caption{\label{fig-regCC1mod} Original version. Using model controlling for age, age$^2$, male,
married, employment type dummies(full/part time, self-employed, retired,
housewife, student, unemployed, and other), financial satisfaction, health, kids,
and religion important.  See appendix for results in table also including
omitted here $20-500k$ category. 
See appendix model for
pooled sample in table \ref{regB} except that
original 8 step urbanicity variable is now categorized into 3 bins as there are
smaller cell sizes separately by country.
%
 Still note that few countries lack the estimates as even
this simplified 3 step urbanicity spectrum is not available. 
Overall the results are mostly insignificant,
but more negative (11) than positive (6). So in general within countries for few
significant relationships there is less freedom in
smaller areas as theorized in the literature.}
\end{figure}


Only 16 out of 64 countries have more freedom in cities. 
They have at least one coefficient on either $<20k$ or $20-500k$ (v
$>500k$) negative and significant at $<0.05$. 6 countries, on the other hand,
have more freedom in ruralities--they have positive effect coefficients on
smaller areas dummies. Note that figure \ref{fig-regCC1mod} only shows $<20k$
coefficients; for $20-500k$ coefficients see appendix.


It is difficult to see any consistent country patterns in the table. Perhaps
what stands out is that among the richer or more economically developed
countries there are smaller differences.
 %
In poorer/less economically
developed countries there are more differences: mostly less freedom in smaller
places.\footnote{But sometimes more.} This makes sense--smaller places in
developing countries can be very traditional and very poor--freedom is difficult
in such places. In richer countries, on the other hand, one can already be free
just about anywhere, no need for a large city. This is similar to what happiness
literature have found with respect to happiness \citep{aokCityBook15}--only in
the very poorest places of all such as Sub-Saharan Africa happiness is greater in
larger cities as life is unbearable and lacking necessities such as shelter and
clean water in ruralities.



Countries could by classified by region (see appendix).
About half of Latin countries in this sample have lower freedom in smaller areas: ARG, BOL, CHL, NIC, and
PER; and about the other half are insignificant: BRA, COL, ECU, GTM, PRI, URY, and VEN.
 
South Asia does not have a clear pattern: BGD is positive and PAK marginally positive, but IND highly
negative, and MDV is insignificant.
 
Europe/Central Asia: mostly insignificant: ARM, CZE, GRC, KGZ, NLD, RUS, SRB, SVK,
and UKR. But otherwise mostly negative: AND marginally, CYP, KAZ, TJK marginally, and TUR;
and positive: ROU, DEU marginally. Note that GBR has missing obs and does not
appear in the table at all, and UZB has missing urbanicity data. 

East Asia/Pacific is mostly insignificant: AUS, CHN, IDN, MMR, NZL, PHL.  
%East Asia/Pacific stood out earlier as mostly positive, but here: negative
And several countries in East Asia/Pacific are without respondents in smaller
areas: HKG, KOR, MAC, SGP, and VNM.
Positive: MNG, and highly positive THA; 
and negative: JPN, MYS.

Middle East/North Africa mostly insignificant: EGY, IRN, JOR, LBY,  and TUN.  
 Two negative: IRQ, LBN; and one positive: MAR.

 Sub-Saharan Africa: two  insignificant: KEN, NGA; ETH is negative, but ZWE positive.

 Two countries missing the region classification by the World Bank: NIR and TWN.
 NIR has a small sample size of less than 500 and missing observations on the variables
 used in regression so it does not appear in the table. TWN is insignificant. 



\subsection{Exploring political say}

\begin{figure}[H]
 \includegraphics[height=3in]{m-body.pdf}\centering
\caption{\label{m-body} Multilevel country random intercept regressions of
  freedom on rural-urban dummies interacted with: a) financial satisfaction, and b)
  class. For underlying regressions see appendix. Controls are
  the same as in the previous section.}
\end{figure}

 \begin{figure}[H]
  \includegraphics[height=3in]{m-polSay.pdf}\centering
  \caption{\label{m-polSay} mixed polSay $i.t4\#\#i.sf3$    age age2 male mar
    i.emp health kids $rel\_imp$  $\|\|$ cc: and the second one same just
    interacted with soc class}
 \end{figure}


GDP took avg for 2015-2025. And same results as for the World Bank income
groups. \\

Political say variable--the sample size drops quite a lot--overall there are
about 94k observations. But this variable has 69k only; with freedom have more
like 85k

\begin{scriptsize}
\begin{spacing}{.19} 
\begin{verbatim}
mixed polSay i.t3##c.ny_gdp_pcap_kd  mar i.emp satFin health kids rel_imp  || cc: i.t3  ,mle

 ta cc t3 if e(sample)==1

ISO 3166-1 |
   alpha-3 |
   country |   RECODE of town (place size)
      code |     lt20k    20-500k     gt500k |     Total
-----------+---------------------------------+----------
       AND |       718        270          0 |       988 
       ARG |       124        135        467 |       726 
       ARM |       497        211        418 |     1,126 
       AUS |       412        514        656 |     1,582 
       BGD |       942        201          0 |     1,143 
       BRA |       164        778        640 |     1,582 
       CAN |     1,098      1,694      1,226 |     4,018 
       COL |       304        722        494 |     1,520 
       CYP |       528        358          0 |       886 
       CZE |       668        360        148 |     1,176 
       EGY |       336        388        261 |       985 
       ETH |       812        371          9 |     1,192 
       GTM |        67        450        590 |     1,107 
       HKG |         0          0      2,040 |     2,040 
       IDN |     2,631        331         89 |     3,051 
       IRN |       388        571        472 |     1,431 
       IRQ |       305        219        625 |     1,149 
       JOR |       363        287        490 |     1,140 
       JPN |        43        613        403 |     1,059 
       KAZ |       367        373        316 |     1,056 
       KEN |       479        413        116 |     1,008 
       KGZ |       666        244        193 |     1,103 
       LBN |       407        657        130 |     1,194 
       LBY |       397        409        112 |       918 
       MAC |         0          0        986 |       986 
       MAR |       380        480        340 |     1,200 
       MDV |       705        294          0 |       999 
       MEX |       801        435        384 |     1,620 
       MMR |     1,116         84          0 |     1,200 
       MNG |       378        257      1,003 |     1,638 
       MYS |       744        428        141 |     1,313 
       NGA |       798        352         14 |     1,164 
       NIC |       381        709        110 |     1,200 
       NLD |        81      1,350        162 |     1,593 
       NZL |       363        333        164 |       860 
       PAK |     1,185        230        439 |     1,854 
       PER |       336        928        100 |     1,364 
       PHL |       908        290          0 |     1,198 
       SGP |         0          0      1,940 |     1,940 
       SVK |       723        435          0 |     1,158 
       TJK |       880        210        110 |     1,200 
       TUN |       433        529         77 |     1,039 
       UKR |       535        301        250 |     1,086 
       URY |       255        254        426 |       935 
       VEN |       179        659        349 |     1,187 
       VNM |         0      1,049        151 |     1,200 
       ZWE |       915        265          0 |     1,180 
-----------+---------------------------------+----------
     Total |    24,812     20,441     17,041 |    62,294 
\end{verbatim}
\end{spacing}{.19}
\end{scriptsize}

\begin{scriptsize}
\begin{spacing}{.19} 
\begin{verbatim}


mixed free i.t3##c.ny_gdp_pcap_kd  mar i.emp satFin health kids rel_imp  || cc:  i.t3 ,mle
ta cc t3 if e(sample)==1

ISO 3166-1 |
   alpha-3 |
   country |   RECODE of town (place size)
      code |     lt20k    20-500k     gt500k |     Total
-----------+---------------------------------+----------
       AND |       725        272          0 |       997 
       ARG |       145        158        546 |       849 
       ARM |       532        226        434 |     1,192 
       AUS |       411        515        658 |     1,584 
       BGD |       988        209          0 |     1,197 
       BOL |       590        711        689 |     1,990 
       BRA |       169        826        693 |     1,688 
       CAN |     1,098      1,694      1,226 |     4,018 
       CHL |       141        436        326 |       903 
       CHN |        48      2,727          0 |     2,775 
       COL |       304        722        494 |     1,520 
       CYP |       571        372          0 |       943 
       CZE |       675        361        150 |     1,186 
       DEU |       693        604        194 |     1,491 
       ECU |       375        462        347 |     1,184 
       EGY |       416        449        310 |     1,175 
       ETH |       830        376          9 |     1,215 
       GRC |       406        324        431 |     1,161 
       GTM |        67        454        592 |     1,113 
       HKG |         0          0      2,060 |     2,060 
       IDN |     2,722        341         91 |     3,154 
       IRN |       396        598        483 |     1,477 
       IRQ |       311        228        638 |     1,177 
       JOR |       375        305        522 |     1,202 
       JPN |        45        654        431 |     1,130 
       KAZ |       385        399        332 |     1,116 
       KEN |       479        413        119 |     1,011 
       KGZ |       678        258        195 |     1,131 
       KOR |         0        364        881 |     1,245 
       LBN |       410        660        130 |     1,200 
       LBY |       403        412        115 |       930 
       MAC |         0          0        984 |       984 
       MAR |       380        480        340 |     1,200 
       MDV |       710        300          0 |     1,010 
       MEX |       802        435        384 |     1,621 
       MMR |     1,116         84          0 |     1,200 
       MNG |       378        257      1,003 |     1,638 
       MYS |       744        428        141 |     1,313 
       NGA |       829        369         14 |     1,212 
       NIC |       381        709        110 |     1,200 
       NLD |        86      1,416        173 |     1,675 
       NZL |       361        329        166 |       856 
       PAK |     1,261        243        471 |     1,975 
       PER |       347        940        100 |     1,387 
       PHL |       909        290          0 |     1,199 
       PRI |        72        993          0 |     1,065 
       ROU |       630        368        101 |     1,099 
       RUS |       539        577        560 |     1,676 
       SGP |         0          0      1,966 |     1,966 
       SRB |       372        427        184 |       983 
       SVK |       735        440          0 |     1,175 
       THA |     1,064         96        107 |     1,267 
       TJK |       880        210        110 |     1,200 
       TUN |       448        548         77 |     1,073 
       TUR |       232      1,724        402 |     2,358 
       UKR |       553        309        254 |     1,116 
       URY |       258        272        436 |       966 
       USA |         0      1,395        363 |     1,758 
       VEN |       179        659        349 |     1,187 
       VNM |         0      1,049        151 |     1,200 
       ZWE |       933        269          0 |     1,202 
-----------+---------------------------------+----------
     Total |    30,587     32,146     22,042 |    84,775 


\end{verbatim}
\end{spacing}{.19}
\end{scriptsize}



 \begin{figure}[H]
  \includegraphics[height=3in]{original.pdf}\centering
  \caption{\label{original} reg polSay $i.t3\#\#i.iii$ mar $i.emp$ satFin health kids $rel\_imp$ , $robust cluster(ccc)$}
 \end{figure}
 \begin{figure}[H]
  \includegraphics[height=3in]{originalD.pdf}\centering
  \caption{\label{originalD} reg polSay $i.t3\#\#i.iii$ mar $i.emp$ satFin health kids $rel\_imp$ , $robust cluster(ccc)$}
 \end{figure}


 \begin{figure}[H]
  \includegraphics[height=3in]{m-ps1.pdf}\centering
  \caption{\label{m-ps1} mixed polSay $i.t3\#\#i.iii$  mar i.emp satFin health kids rel\_imp $ \|\|$ cc: i.t3  ,mle}
 \end{figure}
 \begin{figure}[H]
  \includegraphics[height=3in]{m-ps1D.pdf}\centering
  \caption{\label{m-ps1D} mixed polSay $i.t3\#\#i.iii$  mar i.emp satFin health kids rel\_imp $ \|\|$ cc: i.t3  ,mle}
 \end{figure}

\begin{figure}[H]
 \includegraphics[height=3in]{m-psGdp.pdf}\centering
 \caption{\label{m-psGdp} }
\end{figure}


 
 \begin{figure}[H]
  \includegraphics[height=3in]{m-f1.pdf}\centering
  \caption{\label{m-f1} mixed free $i.t3\#\#i.iii$  mar i.emp satFin health kids rel\_imp  $\|\|$ cc:  i.t3 ,mle}
 \end{figure}
 \begin{figure}[H]
  \includegraphics[height=3in]{m-f1D.pdf}\centering
  \caption{\label{m-f1D} mixed free $i.t3\#\#i.iii$  mar i.emp satFin health kids rel\_imp  $\|\|$ cc:  i.t3 ,mle}
 \end{figure}

  \begin{figure}[H]
  \includegraphics[height=3in]{m-fGdp.pdf}\centering
  \caption{\label{m-fGdp} }
 \end{figure}

  \begin{figure}[H]
  \includegraphics[height=3in]{m-fGdpB.pdf}\centering
  \caption{\label{m-fGdpB} without controls freedom is not lowest for richest places}
 \end{figure}
 
 
 
 \begin{figure}[H]
  \includegraphics[height=3in]{m-polSay.pdf}\centering
  \caption{\label{m-polSay} mixed polSay $i.t4\#\#i.sf3$    age age2 male mar
    i.emp health kids $rel\_imp$  $\|\|$ cc: and the second one same just
    interacted with soc class}
 \end{figure}
 


 \begin{figure}[H]
  \includegraphics[height=3in]{psBar.pdf}\centering
  \caption{\label{psBar} }
 \end{figure}



\begin{scriptsize}
\begin{spacing}{.19} 
\begin{verbatim}
. reg polSay TTT2 TTT3  mar i.emp satFin health kids rel_imp i.ccc if iii==1, robust

Linear regression                               Number of obs     =     20,590
                                                F(30, 20559)      =     114.51
                                                Prob > F          =     0.0000
                                                R-squared         =     0.1176
                                                Root MSE          =     1.2078

--------------------------------------------------------------------------------
               |               Robust
        polSay | Coefficient  std. err.      t    P>|t|     [95% conf. interval]
---------------+----------------------------------------------------------------
          TTT2 |      -0.04       0.02    -1.92    0.05        -0.08        0.00
          TTT3 |      -0.06       0.03    -2.24    0.02        -0.12       -0.01
           mar |       0.07       0.02     3.25    0.00         0.03        0.12
               |
           emp |
    Part time  |      -0.06       0.03    -1.70    0.09        -0.12        0.01
Self employed  |      -0.04       0.03    -1.74    0.08        -0.09        0.01
      Retired  |      -0.02       0.04    -0.58    0.56        -0.10        0.05
    Housewife  |      -0.13       0.03    -4.80    0.00        -0.19       -0.08
     Students  |      -0.01       0.04    -0.32    0.75        -0.09        0.06
   Unemployed  |      -0.07       0.03    -2.25    0.02        -0.14       -0.01
        Other  |       0.06       0.13     0.43    0.67        -0.21        0.32
               |
        satFin |       0.03       0.00     7.74    0.00         0.02        0.04
        health |       0.08       0.01     7.95    0.00         0.06        0.10
          kids |      -0.00       0.01    -0.27    0.79        -0.01        0.01
       rel_imp |       0.03       0.01     2.10    0.04         0.00        0.06
               |
           ccc |
          104  |      -0.72       0.05   -14.37    0.00        -0.82       -0.62
          231  |      -0.11       0.05    -2.29    0.02        -0.21       -0.02
          404  |      -0.54       0.05   -10.13    0.00        -0.64       -0.43
          417  |      -0.58       0.05   -10.91    0.00        -0.68       -0.48
          496  |      -0.88       0.05   -17.23    0.00        -0.98       -0.78
          504  |      -1.11       0.04   -24.79    0.00        -1.20       -1.02
          558  |      -0.51       0.06    -9.21    0.00        -0.62       -0.40
          566  |      -0.82       0.06   -14.69    0.00        -0.93       -0.71
          586  |      -0.42       0.05    -9.29    0.00        -0.51       -0.33
          608  |       0.28       0.04     6.47    0.00         0.20        0.36
          704  |      -0.09       0.05    -1.67    0.09        -0.20        0.02
          716  |      -0.64       0.05   -11.99    0.00        -0.75       -0.54
          762  |      -0.08       0.05    -1.76    0.08        -0.18        0.01
          788  |      -0.48       0.05    -9.71    0.00        -0.58       -0.38
          804  |      -1.17       0.05   -21.77    0.00        -1.28       -1.07
          818  |      -0.97       0.05   -20.06    0.00        -1.06       -0.87
               |
         _cons |       2.91       0.08    36.31    0.00         2.76        3.07
--------------------------------------------------------------------------------

. 
end of do-file

. do "/tmp/SD18984.000000"

. reg polSay TTT2 TTT3  mar i.emp satFin health kids rel_imp i.ccc if iii==3, robust

Linear regression                               Number of obs     =     19,221
                                                F(26, 19194)      =      82.06
                                                Prob > F          =     0.0000
                                                R-squared         =     0.1055
                                                Root MSE          =     1.0325

--------------------------------------------------------------------------------
               |               Robust
        polSay | Coefficient  std. err.      t    P>|t|     [95% conf. interval]
---------------+----------------------------------------------------------------
          TTT2 |       0.13       0.02     5.93    0.00         0.09        0.18
          TTT3 |       0.21       0.03     7.97    0.00         0.16        0.26
           mar |       0.08       0.02     4.89    0.00         0.05        0.12
               |
           emp |
    Part time  |      -0.03       0.03    -1.26    0.21        -0.09        0.02
Self employed  |      -0.07       0.03    -2.11    0.03        -0.14       -0.01
      Retired  |      -0.04       0.02    -1.96    0.05        -0.08       -0.00
    Housewife  |      -0.13       0.03    -3.65    0.00        -0.19       -0.06
     Students  |       0.11       0.03     3.25    0.00         0.05        0.18
   Unemployed  |      -0.11       0.04    -2.63    0.01        -0.19       -0.03
        Other  |      -0.13       0.06    -2.16    0.03        -0.26       -0.01
               |
        satFin |       0.06       0.00    14.19    0.00         0.05        0.07
        health |       0.06       0.01     5.45    0.00         0.04        0.08
          kids |      -0.01       0.01    -1.72    0.09        -0.03        0.00
       rel_imp |       0.08       0.01    10.77    0.00         0.07        0.10
               |
           ccc |
           36  |      -0.62       0.05   -12.64    0.00        -0.72       -0.52
          124  |      -0.35       0.04    -7.73    0.00        -0.43       -0.26
          196  |      -0.86       0.06   -15.29    0.00        -0.97       -0.75
          203  |      -1.00       0.05   -19.61    0.00        -1.10       -0.90
          344  |      -0.55       0.05   -10.76    0.00        -0.65       -0.45
          392  |       0.09       0.06     1.50    0.13        -0.03        0.20
          446  |      -0.35       0.05    -6.33    0.00        -0.46       -0.24
          528  |      -0.61       0.05   -12.20    0.00        -0.71       -0.52
          554  |      -0.49       0.05    -9.64    0.00        -0.59       -0.39
          702  |      -0.52       0.05    -9.86    0.00        -0.63       -0.42
          703  |      -1.06       0.05   -20.09    0.00        -1.16       -0.95
          858  |      -0.50       0.06    -8.72    0.00        -0.61       -0.39
               |
         _cons |       2.46       0.06    40.03    0.00         2.34        2.58
--------------------------------------------------------------------------------

. 
end of do-file


. reg polSay i.t3##i.iii  mar i.emp satFin health kids rel_imp  i.ccc, robust
note: 818.ccc omitted because of collinearity.
note: 862.ccc omitted because of collinearity.

Linear regression                               Number of obs     =     62,294
                                                F(64, 62229)      =     125.62
                                                Prob > F          =     0.0000
                                                R-squared         =     0.1008
                                                Root MSE          =     1.1565

--------------------------------------------------------------------------------
               |               Robust
        polSay | Coefficient  std. err.      t    P>|t|     [95% conf. interval]
---------------+----------------------------------------------------------------
            t3 |
      20-500k  |      -0.04       0.02    -1.83    0.07        -0.08        0.00
       gt500k  |      -0.06       0.03    -2.14    0.03        -0.11       -0.01
               |
           iii |
            2  |       0.24       0.05     4.60    0.00         0.14        0.35
            3  |       0.79       0.06    14.20    0.00         0.68        0.90
               |
        t3#iii |
    20-500k#2  |      -0.04       0.03    -1.32    0.19        -0.10        0.02
    20-500k#3  |       0.18       0.03     5.67    0.00         0.11        0.24
     gt500k#2  |      -0.04       0.04    -1.22    0.22        -0.12        0.03
     gt500k#3  |       0.27       0.04     7.11    0.00         0.20        0.35
               |
           mar |       0.06       0.01     5.61    0.00         0.04        0.09
               |
           emp |
    Part time  |      -0.00       0.02    -0.12    0.91        -0.04        0.03
Self employed  |      -0.04       0.02    -2.33    0.02        -0.07       -0.01
      Retired  |      -0.02       0.02    -0.93    0.35        -0.05        0.02
    Housewife  |      -0.09       0.02    -5.51    0.00        -0.12       -0.06
     Students  |       0.04       0.02     1.91    0.06        -0.00        0.08
   Unemployed  |      -0.05       0.02    -2.61    0.01        -0.09       -0.01
        Other  |      -0.04       0.05    -0.69    0.49        -0.14        0.07
               |
        satFin |       0.03       0.00    15.08    0.00         0.03        0.04
        health |       0.07       0.01    11.71    0.00         0.06        0.08
          kids |       0.00       0.00     0.13    0.90        -0.01        0.01
       rel_imp |       0.06       0.01    11.16    0.00         0.05        0.08
               |
           ccc |
           32  |       1.21       0.05    26.33    0.00         1.12        1.30
           36  |      -0.63       0.05   -12.97    0.00        -0.73       -0.54
           50  |       0.96       0.05    20.24    0.00         0.87        1.05
           51  |       0.18       0.05     3.47    0.00         0.08        0.29
           76  |      -0.30       0.05    -6.38    0.00        -0.39       -0.21
          104  |       0.26       0.05     4.93    0.00         0.15        0.36
          124  |      -0.35       0.04    -7.74    0.00        -0.43       -0.26
          170  |       0.01       0.05     0.20    0.84        -0.09        0.11
          196  |      -0.86       0.06   -15.39    0.00        -0.97       -0.75
          203  |      -1.00       0.05   -19.83    0.00        -1.10       -0.91
          231  |       0.87       0.05    17.58    0.00         0.77        0.97
          320  |       0.13       0.05     2.51    0.01         0.03        0.23
          344  |      -0.55       0.05   -10.75    0.00        -0.65       -0.45
          360  |       0.60       0.04    13.81    0.00         0.52        0.69
          364  |       0.22       0.05     4.65    0.00         0.13        0.32
          368  |       0.09       0.05     1.69    0.09        -0.01        0.19
          392  |       0.06       0.06     0.99    0.32        -0.05        0.17
          398  |       0.65       0.05    13.65    0.00         0.56        0.75
          400  |       0.48       0.05     9.43    0.00         0.38        0.58
          404  |       0.45       0.05     8.56    0.00         0.34        0.55
          417  |       0.41       0.05     7.87    0.00         0.30        0.51
          422  |      -0.03       0.05    -0.61    0.54        -0.13        0.07
          434  |       0.07       0.05     1.21    0.23        -0.04        0.17
          446  |      -0.35       0.05    -6.34    0.00        -0.45       -0.24
          458  |       0.46       0.05     9.66    0.00         0.37        0.56
          462  |      -0.28       0.05    -5.29    0.00        -0.38       -0.18
          484  |       0.02       0.04     0.53    0.59        -0.06        0.11
          496  |       0.16       0.04     3.59    0.00         0.07        0.24
          504  |      -0.14       0.04    -3.20    0.00        -0.23       -0.05
          528  |      -0.61       0.05   -12.19    0.00        -0.71       -0.51
          554  |      -0.50       0.05    -9.84    0.00        -0.60       -0.40
          558  |       0.47       0.05     8.65    0.00         0.37        0.58
          566  |       0.16       0.06     2.86    0.00         0.05        0.27
          586  |       0.55       0.05    11.79    0.00         0.46        0.64
          604  |       0.27       0.05     5.73    0.00         0.18        0.36
          608  |       1.25       0.04    27.97    0.00         1.16        1.34
          702  |      -0.52       0.05    -9.79    0.00        -0.62       -0.41
          703  |      -1.07       0.05   -20.44    0.00        -1.17       -0.97
          704  |       0.94       0.05    18.73    0.00         0.84        1.03
          716  |       0.34       0.05     6.32    0.00         0.23        0.44
          762  |       0.90       0.05    18.66    0.00         0.80        0.99
          788  |       0.49       0.05    10.10    0.00         0.40        0.59
          804  |      -0.15       0.05    -3.02    0.00        -0.25       -0.05
          818  |       0.00  (omitted)
          858  |      -0.51       0.06    -8.96    0.00        -0.62       -0.40
          862  |       0.00  (omitted)
               |
         _cons |       1.80       0.05    36.71    0.00         1.71        1.90
--------------------------------------------------------------------------------

. 

\end{verbatim}
\end{spacing}{.19}
\end{scriptsize}


 \begin{figure}[H]
  \includegraphics[height=3in]{m-ps1.pdf}\centering
  \caption{\label{m-ps1} por some reason margins dont work for me with cc
    dummies this is without cc dummies, but above in stata printout interactions
  significant with cc dummies}
 \end{figure}















\section{Additional info}

\subsection{AI literature review}
Gemini query ``scholarly literature on individual autonomy or freedom across
urban-rural''\footnote{\url{https://www.google.com/search?q=scholarly+literature+on+individual+autonomy+or+freedom+across+urban-rural&rlz=1CAHWSA_enUS1125US1125&oq=scholarly+literature+on+individual+autonomy+or+freedom+across+urban-rural&gs_lcrp=EgZjaHJvbWUyBggAEEUYOdIBCTIyNDYyajBqN6gCCLACAfEFAI5zaHyURfbxBQCOc2h8lEX2&sourceid=chrome&source=chrome.ob&ie=UTF-8}}

\begin{table}[!ht]
  \centering\scriptsize

  
    \begin{tabular}{|l|l|l|l|l|l|l|l|l|l|}
    \hline
        Autonomy Dimension & Rural & Urban\\ \hline
        Social Control & Strict; governed by kinship, mutual surveillance, and traditional norms. & Loose; regulated by formal systems (laws, courts) and shielded by anonymity.\\ \hline
        Core Values & Interdependence, community solidarity, and generational stability. & Individualism, open-mindedness, and socioeconomic mobility.\\ \hline
        Family Dynamics & High value placed on obedience and material/financial cooperation. & Focus on psychological independence and diverse career/lifestyle paths.\\ \hline
        Physical Liberty & High geographic isolation; dependency on personal transit for survival. & High physical access via transit; restricted by crowd density and safety rules.\\ \hline
    \end{tabular}
\end{table}

%suggested the following literature summarised in this section by us, the human authors:
%put into body lord2023autonomy pospvech2024loss diamond2023 mishra2014relative


\subsection{Country sample sizes; Regional country classification; Variables'
  definitions (survey questions); Histograms}

We use World Bank regions classification into several regions:
 East Asia/Pacific, 
 Europe/Central Asia, 
 Latin America/Caribbean, 
 Middle East/North Africa, 
 North America, 
 South Asia, and
 Sub-Saharan Africa. Countries in each region and sample sizes are in appendix.


\begin{scriptsize}
\begin{spacing}{.19} 
\begin{verbatim}
. bys region: ta cc

------------------------------------------------------------------------------
-> regionname = 

 ISO 3166-1 |
    alpha-3 |
    country |
       code |      Freq.     Percent        Cum.
------------+-----------------------------------
        NIR |        447       26.77       26.77
        TWN |      1,223       73.23      100.00
------------+-----------------------------------
      Total |      1,670      100.00

------------------------------------------------------------------------------
-> regionname = East Asia/Pacific

 ISO 3166-1 |
    alpha-3 |
    country |
       code |      Freq.     Percent        Cum.
------------+-----------------------------------
        AUS |      1,813        7.29        7.29
        CHN |      3,036       12.21       19.50
        HKG |      2,075        8.35       27.85
        IDN |      3,200       12.87       40.72
        JPN |      1,353        5.44       46.16
        KOR |      1,245        5.01       51.16
        MAC |      1,023        4.11       55.28
        MMR |      1,200        4.83       60.10
        MNG |      1,638        6.59       66.69
        MYS |      1,313        5.28       71.97
        NZL |      1,057        4.25       76.22
        PHL |      1,200        4.83       81.05
        SGP |      2,012        8.09       89.14
        THA |      1,500        6.03       95.17
        VNM |      1,200        4.83      100.00
------------+-----------------------------------
      Total |     24,865      100.00

------------------------------------------------------------------------------
-> regionname = Europe/Central Asia

 ISO 3166-1 |
    alpha-3 |
    country |
       code |      Freq.     Percent        Cum.
------------+-----------------------------------
        AND |      1,004        3.88        3.88
        ARM |      1,223        4.73        8.61
        CYP |      1,000        3.87       12.48
        CZE |      1,200        4.64       17.12
        DEU |      1,528        5.91       23.03
        GBR |      2,609       10.09       33.13
        GRC |      1,200        4.64       37.77
        KAZ |      1,276        4.94       42.70
        KGZ |      1,200        4.64       47.35
        NLD |      2,145        8.30       55.64
        ROU |      1,257        4.86       60.51
        RUS |      1,810        7.00       67.51
        SRB |      1,046        4.05       71.55
        SVK |      1,200        4.64       76.20
        TJK |      1,200        4.64       80.84
        TUR |      2,415        9.34       90.18
        UKR |      1,289        4.99       95.16
        UZB |      1,250        4.84      100.00
------------+-----------------------------------
      Total |     25,852      100.00

------------------------------------------------------------------------------
-> regionname = Latin America/Caribbean

 ISO 3166-1 |
    alpha-3 |
    country |
       code |      Freq.     Percent        Cum.
------------+-----------------------------------
        ARG |      1,003        5.75        5.75
        BOL |      2,067       11.85       17.60
        BRA |      1,762       10.10       27.71
        CHL |      1,000        5.73       33.44
        COL |      1,520        8.72       42.16
        ECU |      1,200        6.88       49.04
        GTM |      1,229        7.05       56.09
        MEX |      1,741        9.98       66.07
        NIC |      1,200        6.88       72.95
        PER |      1,400        8.03       80.98
        PRI |      1,127        6.46       87.44
        URY |      1,000        5.73       93.18
        VEN |      1,190        6.82      100.00
------------+-----------------------------------
      Total |     17,439      100.00

------------------------------------------------------------------------------
-> regionname = Middle East/North Africa

 ISO 3166-1 |
    alpha-3 |
    country |
       code |      Freq.     Percent        Cum.
------------+-----------------------------------
        EGY |      1,200       12.11       12.11
        IRN |      1,499       15.13       27.25
        IRQ |      1,200       12.11       39.36
        JOR |      1,203       12.14       51.50
        LBN |      1,200       12.11       63.62
        LBY |      1,196       12.07       75.69
        MAR |      1,200       12.11       87.81
        TUN |      1,208       12.19      100.00
------------+-----------------------------------
      Total |      9,906      100.00

------------------------------------------------------------------------------
-> regionname = North America

 ISO 3166-1 |
    alpha-3 |
    country |
       code |      Freq.     Percent        Cum.
------------+-----------------------------------
        CAN |      4,018       60.75       60.75
        USA |      2,596       39.25      100.00
------------+-----------------------------------
      Total |      6,614      100.00

------------------------------------------------------------------------------
-> regionname = South Asia

 ISO 3166-1 |
    alpha-3 |
    country |
       code |      Freq.     Percent        Cum.
------------+-----------------------------------
        BGD |      1,200       20.25       20.25
        IND |      1,692       28.55       48.80
        MDV |      1,039       17.53       66.33
        PAK |      1,995       33.67      100.00
------------+-----------------------------------
      Total |      5,926      100.00

------------------------------------------------------------------------------
-> regionname = Sub-Saharan Africa

 ISO 3166-1 |
    alpha-3 |
    country |
       code |      Freq.     Percent        Cum.
------------+-----------------------------------
        ETH |      1,230       24.86       24.86
        KEN |      1,266       25.59       50.44
        NGA |      1,237       25.00       75.44
        ZWE |      1,215       24.56      100.00
------------+-----------------------------------
      Total |      4,948      100.00

\end{verbatim}
\end{spacing}{.19}
\end{scriptsize}


\input{varDes.tex}

Per distributions it is worth noting left-skewed distribution on freedom,
unexpectedly most people feel maximally free at 10 out of 1-10, likewise but as expected with
urbanism--most people are in the top category $>500k$, and again left-skewed on
religion, probably not unexpected--religiosity is declining, but still
strong. 

\input{hist.tex}

\subsection{First/initial results}
\subsubsection{by region}

These were our first/initial results but they do not hold when including country
dummies and so we postpone them to appendix. 

First for all regions in general there is more freedom in smaller areas v top
bin, the base case ($>500k$) in table \ref{regA} but controlling for other
predictors in table \ref{regB} it is just the 2nd (2-5k) and 3rd (5-10k)
smallest categories, and smaller values in 6th and 7th categories.
But the very smallest areas ($<2k$) are not more free.  
% indeed negative -.07 for the smallest areas ($<2k$).
 Next we separate by region.

East Asia/Pacific stands out with high effect sizes at about .4 or even .5, also
controlling for other predictors of freedom. Europe is mostly less free in
smaller areas. Latin America/Caribbean
stands out with negative effect in the very smallest areas. Middle East/North
Africa has mostly negative effects. North America is positive in 2nd largest
category (100-500k). South Asia and Sub-Saharan Africa mostly do not show significant
differences.\footnote{Except .3 for South Asia for the smallest places ($<$2k). %-0.43 for Sub-Saharan Africa in 7th bin
  % (100-500k).
} 

\begin{spacing}{.9} \begin{table}[H]\centering\begin{scriptsize}
\begin{tabular}{p{.8in}p{.55in}p{.55in}p{.55in}p{.55in}p{.55in}p{.55in}p{.55in}p{.55in}p{.55in}}\hline
                    &   $<.5m$   &     $>.5m$   &   $<.5m$   &     $>.5m$   &   URYrurTow   &     URYcity   \\
2022                &       -0.21** &       -0.41** &       -0.12** &       -0.23   &        0.75***&        0.23+  \\
constant            &        7.54***&        7.65***&        7.50***&        7.38***&        7.54***&        7.69***\\
N                   &        3111   &         521   &        3572   &         373   &        1154   &         836   \\
\hline + 0.10 * 0.05 ** 0.01 *** 0.001; robust std err
\end{tabular}\end{scriptsize}\caption{\label{regA}OLS regressions of freedom/agency}
\end{table}\end{spacing}

\begin{spacing}{.9} \begin{table}[H]\centering\begin{scriptsize}
\begin{tabular}{p{.8in}p{.55in}p{.55in}p{.55in}p{.55in}p{.55in}p{.55in}p{.55in}p{.55in}p{.55in}}\hline
                    &   $<.5m$   &     $>.5m$   &   $<.5m$   &     $>.5m$   &   URYrurTow   &     URYcity   \\
2022                &       -0.18*  &       -0.39+  &       -0.20***&       -0.45** &        0.42***&        0.21   \\
income              &        0.09***&        0.01   &        0.06***&        0.14***&        0.07*  &        0.13***\\
age                 &       -0.03*  &       -0.08** &       -0.02+  &       -0.06+  &        0.00   &       -0.06** \\
age2                &        0.00** &        0.00** &        0.00** &        0.00*  &       -0.00   &        0.00** \\
male                &       -0.18** &       -0.13   &       -0.11*  &       -0.27+  &        0.06   &        0.19   \\
married or living together as married&        0.53***&        0.74***&        0.44***&        0.23   &        0.46** &        0.06   \\
divorced/separated/widowed&        0.07   &        0.15   &       -0.11   &       -0.14   &       -0.37+  &       -0.19   \\
autonomy            &       -0.11*  &       -0.07   &       -0.11** &       -0.01   &       -0.06   &        0.06   \\
freedom             &        0.44***&        0.42***&        0.35***&        0.43***&        0.43***&        0.36***\\
trust               &        0.12+  &        0.42** &        0.43***&        0.28+  &       -0.05   &        0.10   \\
postmaterialist     &       -0.05   &       -0.18   &       -0.11*  &        0.14   &       -0.02   &        0.15   \\
god important       &        0.01   &        0.05*  &        0.02*  &       -0.01   &        0.05** &        0.06** \\
constant            &        4.08***&        5.95***&        4.59***&        4.80***&        3.47***&        4.58***\\
N                   &        1985   &         309   &        2283   &         237   &         736   &         579   \\
\hline + 0.10 * 0.05 ** 0.01 *** 0.001; robust std err
\end{tabular}\end{scriptsize}\caption{\label{regB}OLS regressions of freedom/agency}
\end{table}\end{spacing}

\subsubsection{Insignificant results with country dummies}

Country dummies render cross-country patters insignificant. 

Here we conduct important sensitivity/robustness checks. First in a1 just
urbanicity dummies. Then a1a adds region dummies like earlier in the first table
\ref{regA} little change. A1b has income level dummies TODO spell out; again
little change; but addition of country dummies in a1c kills the significance and
turns many categories negative notably the first smallest one, similar with
controls in models a2* and country dummies make categories even more negative
in a2c than in a1c. And then a1c2 and a2c2 add clustered standard errors on
country (stata: robust cluster(country))--all significance disappears except in
a2c2 on the smallest areas--hence only the very smallest ares $<2k$ are less free
than the largest ones $>500k$.



in table \ref{regAA}  we conduct additional robustness checks, a1 and a2 are
repeated from the body; a1a and a2a include region dummies: in a1a  $100-500k$ results are
attenuated, but others are stronger; and in a2a the top three categories
spanning $20-500k$ are weaker but other categories are stronger and the smaller
category $<2k$ positive effect goes away. But then in a1c and a2c addition of
country dummies kills the effects unfortunately. And having clustered standard
errors yields all results insignificant


\begin{spacing}{.9} \begin{table}[H]\centering\begin{scriptsize}
\begin{tabular}{p{.8in}p{.55in}p{.55in}p{.55in}p{.55in}p{.55in}p{.55in}p{.55in}p{.55in}p{.55in}p{.55in}p{.55in}}\hline
\input{regAA.tex}\hline + 0.10 * 0.05 ** 0.01 *** 0.001; robust std err
\end{tabular}\end{scriptsize}\caption{\label{regAA}OLS regressions of freedom/agency}
\end{table}\end{spacing}






\subsection{Freedom main results from paper}
\subsubsection{By country}

First let's start with sample sizes by country by 3 urbanicity categories--in
general they are sizable $>100$, and in cases where they are not CI is wider on
the estimates such as with ETH at -1.12 on lt20k that is insignificant

\begin{scriptsize}
\begin{spacing}{.19} 
\begin{verbatim}

ISO 3166-1 |
   alpha-3 |
   country |   RECODE of town (place size)
      code |     lt20k    20-500k     gt500k |     Total
-----------+---------------------------------+----------
       AND |       730        274          0 |     1,004 
       ARG |       200        180        623 |     1,003 
       ARM |       545        232        446 |     1,223 
       AUS |       453        567        705 |     1,725 
       BGD |       991        209          0 |     1,200 
       BOL |       625        741        701 |     2,067 
       BRA |       185        864        713 |     1,762 
       CAN |     1,098      1,694      1,226 |     4,018 
       CHL |       152        487        361 |     1,000 
       CHN |        50      2,986          0 |     3,036 
       COL |       304        722        494 |     1,520 
       CYP |       615        385          0 |     1,000 
       CZE |       680        370        150 |     1,200 
       DEU |       712        616        200 |     1,528 
       ECU |       384        466        350 |     1,200 
       EGY |       424        456        320 |     1,200 
       ETH |       841        380          9 |     1,230 
       GBR |       495      1,713        397 |     2,605 
       GRC |       420        340        440 |     1,200 
       GTM |        71        505        653 |     1,229 
       HKG |         0          0      2,075 |     2,075 
       IDN |     2,761        346         93 |     3,200 
       IND |     1,415        235         42 |     1,692 
       IRN |       400        609        490 |     1,499 
       IRQ |       320        232        648 |     1,200 
       JOR |       375        306        522 |     1,203 
       JPN |        51        786        516 |     1,353 
       KAZ |       457        456        363 |     1,276 
       KEN |       499        421        122 |     1,042 
       KGZ |       710        290        200 |     1,200 
       KOR |         0        364        881 |     1,245 
       LBN |       410        660        130 |     1,200 
       LBY |       480        544        172 |     1,196 
       MAC |         0          0      1,023 |     1,023 
       MAR |       380        480        340 |     1,200 
       MDV |       738        301          0 |     1,039 
       MEX |       858        463        403 |     1,724 
       MMR |     1,116         84          0 |     1,200 
       MNG |       378        257      1,003 |     1,638 
       MYS |       744        428        141 |     1,313 
       NGA |       848        375         14 |     1,237 
       NIC |       381        709        110 |     1,200 
       NIR |        39        407          0 |       446 
       NLD |       104      1,710        217 |     2,031 
       NZL |       420        387        187 |       994 
       PAK |     1,272        249        474 |     1,995 
       PER |       355        945        100 |     1,400 
       PHL |       910        290          0 |     1,200 
       PRI |        73      1,054          0 |     1,127 
       ROU |       703        435        119 |     1,257 
       RUS |       594        626        590 |     1,810 
       SGP |         0          0      2,012 |     2,012 
       SRB |       403        447        196 |     1,046 
       SVK |       750        450          0 |     1,200 
       THA |     1,244        116        140 |     1,500 
       TJK |       880        210        110 |     1,200 
       TUN |       514        614         80 |     1,208 
       TUR |       240      1,767        408 |     2,415 
       TWN |        40      1,165         18 |     1,223 
       UKR |       638        355        296 |     1,289 
       URY |       270        290        440 |     1,000 
       USA |         0      1,443        385 |     1,828 
       UZB |         0          0      1,250 |     1,250 
       VEN |       180        661        349 |     1,190 
       VNM |         0      1,049        151 |     1,200 
       ZWE |       945        270          0 |     1,215 
-----------+---------------------------------+----------
     Total |    33,870     37,473     24,598 |    95,941 
\end{verbatim}
\end{spacing}{.19}
\end{scriptsize}





\begin{spacing}{.9} \begin{table}[H]\centering   \begin{scriptsize}
\begin{tabular}{p{1.8in}p{.5in}p{.5in}p{.5in}p{.5in}p{.5in}p{.5in}p{.5in}}\hline
\input{regCC1mod.tex}
\hline +0.10 *0.05 **0.01 ***0.001; robust std err
\end{tabular}\end{scriptsize}\caption{\label{regCC1mod}
OLS regressions of freedom/autonomy on urbanicity dummies (base: $>500k$) using model controlling for age, age$^2$, male,
married, employment type dummies(full/part time, self-employed, retired,
housewife, student, unemployed, and other), financial satisfaction, health, kids,
and religion important.  See model for pooled sample in table \ref{regB} except that
original 8 step urbanicity variable is now categorized into 3 bins as there are
smaller cell sizes separately by country.
%
 Still note that few countries lack the estimates as even
this simplified 3 step urbanicity spectrum is not available. 
Overall the results are mostly insignificant,
but more negative (11) than positive (6). So in general within countries for few
significant relationships there is less freedom in
smaller areas as theorized in the literature. }
\end{table} \end{spacing}
%regCC1mod: dropped >500k col

Then we split by tertile within a country--the idea is that we will have
retained more observations as oppose to forcing $>500k$ only to be a city, and
we do within country urbanicity where local conditions may be different and
already a place of $100-500k$ may be large. For instance earlier AND had 0
observations in $>500k$ category, but here we have 274 obs in 3rd tertile. 

We were thinking about median split, %leonie: wants to split on median urb-rur
 but we went with tertiles as there are theoretical reasons to contrast rural
 places with large cities \citep{aok-ls_fisher16,aokCityBook15}.

We start below with counts by tertile and country.

\begin{scriptsize}
\begin{spacing}{.19} 
\begin{verbatim}
 ta cc tertile

ISO 3166-1 |
   alpha-3 |
   country |             tertile
      code |         1          2          3 |     Total
-----------+---------------------------------+----------
       AND |       501        229        274 |     1,004 
       ARG |       380        623          0 |     1,003 
       ARM |       448        775          0 |     1,223 
       AUS |       625      1,100          0 |     1,725 
       BGD |       501        335        364 |     1,200 
       BOL |       879      1,188          0 |     2,067 
       BRA |     1,049        713          0 |     1,762 
       CAN |     1,471      1,321      1,226 |     4,018 
       CHL |       639        361          0 |     1,000 
       CHN |     1,980      1,056          0 |     3,036 
       COL |       690        336        494 |     1,520 
       CYP |       367        333        300 |     1,000 
       CZE |       460        400        340 |     1,200 
       DEU |       712        359        457 |     1,528 
       ECU |       509        341        350 |     1,200 
       EGY |       424        456        320 |     1,200 
       ETH |       554        287        389 |     1,230 
       GBR |       989      1,219        397 |     2,605 
       GRC |       420        780          0 |     1,200 
       GTM |       576        653          0 |     1,229 
       HKG |     2,075          0          0 |     2,075 
       IDN |     1,442        791        967 |     3,200 
       IND |       799        389        504 |     1,692 
       IRN |       599        410        490 |     1,499 
       IRQ |       552        648          0 |     1,200 
       JOR |       487        716          0 |     1,203 
       JPN |       837        516          0 |     1,353 
       KAZ |       457        456        363 |     1,276 
       KEN |       395        329        318 |     1,042 
       KGZ |       450        380        370 |     1,200 
       KOR |     1,245          0          0 |     1,245 
       LBN |       410        559        231 |     1,200 
       LBY |       480        349        367 |     1,196 
       MAC |     1,023          0          0 |     1,023 
       MAR |       520        340        340 |     1,200 
       MDV |       494        200        345 |     1,039 
       MEX |       707        614        403 |     1,724 
       MMR |       648        240        312 |     1,200 
       MNG |       575      1,063          0 |     1,638 
       MYS |       580        321        412 |     1,313 
       NGA |       455        393        389 |     1,237 
       NIC |       720        220        260 |     1,200 
       NIR |       161        143        142 |       446 
       NLD |       751      1,063        217 |     2,031 
       NZL |       420        387        187 |       994 
       PAK |       896        438        661 |     1,995 
       PER |       475        825        100 |     1,400 
       PHL |       610        210        380 |     1,200 
       PRI |       726        262        139 |     1,127 
       ROU |       567        343        347 |     1,257 
       RUS |       756        464        590 |     1,810 
       SGP |     2,012          0          0 |     2,012 
       SRB |       365        485        196 |     1,046 
       SVK |       540        390        270 |     1,200 
       THA |       923        138        439 |     1,500 
       TJK |       730         80        390 |     1,200 
       TUN |       442        495        271 |     1,208 
       TUR |     2,007          0        408 |     2,415 
       TWN |       604        601         18 |     1,223 
       UKR |       517        476        296 |     1,289 
       URY |       410        590          0 |     1,000 
       USA |       635        808        385 |     1,828 
       UZB |     1,250          0          0 |     1,250 
       VEN |       401        440        349 |     1,190 
       VNM |     1,049          0        151 |     1,200 
       ZWE |       768        152        295 |     1,215 
-----------+---------------------------------+----------
     Total |    48,139     30,589     17,213 |    95,941 
\end{verbatim}
\end{spacing}{.19}
\end{scriptsize}

Then we repeat twostep graph from the body just withe 1st category as a base in
figure \ref{fff2} and \ref{fff3} so
we can easier compare it to our tertile estimates.

Overall without controls more freedom in cities as hypothesized, but still most
effects are insignificant, and with controls
results are even less significant.

\begin{figure}[H]
 \includegraphics[height=5in]{fff2.pdf}\centering
\caption{\label{fff2} fff2} 
\end{figure}
\begin{figure}[H]
 \includegraphics[height=5in]{fff3.pdf}\centering
\caption{\label{fff3} fff3} 
\end{figure}
 

\begin{figure}[H]
 \includegraphics[height=5in]{l2a.pdf}\centering
\caption{\label{l2a} 2nd tertile: just 2 tertile dummies on rhs, 1st tertile base} 
\end{figure}
\begin{figure}[H]
 \includegraphics[height=5in]{l3a.pdf}\centering
\caption{\label{l3a} 3rd tertile: just 2 tertile dummies on rhs, 1st tertile base} 
\end{figure}

\begin{figure}[H]
 \includegraphics[height=5in]{l2b.pdf}\centering
\caption{\label{l2b} 2nd tertile: all controls} 
\end{figure}
\begin{figure}[H]
 \includegraphics[height=5in]{l3b.pdf}\centering
\caption{\label{l3b} 3rd tertile: all controls} 
\end{figure}





We explore further by interacting region dummies with urban-rural in figure
\ref{t-rr7}. We find significant interactions for Latin America and South Asia
(v East Asia/Pacific)--it is difficult to see patters in figure with all
regions--we plot only the significant interactions in figure \ref{t-rr3}. 

\begin{figure}[H]
 \includegraphics[height=3in]{t-rr7.pdf}\centering
\caption{\label{t-rr7} Interactions of region dummies with urban-rural.}
\end{figure}

\begin{figure}[H]
 \includegraphics[height=3in]{t-rr3.pdf}\centering
\caption{\label{t-rr3} Interactions of region dummies with urban-rural.}
\end{figure}

Last, we interacted urban-rural at person level with \% urban at country level
and the interaction is marginally significant at 0.1 level of significance and
we plot results in figure \ref{t-urb}. Marginsplot shows that in countries that
are less urban, there is more freedom in smaller places, but in more urban
countries there is more freedom in large cities. It is an intriguing finding and
we do not have a good explanation for it. Perhaps in heavily urbanized countries
rural barely exists and it does not help much with freedom, while in more rural
countries there is more meaningful distinction between urban and rural and there
is more freedom in rural areas. 
 

\begin{figure}[H]
 \includegraphics[height=3in]{t-urb.pdf}\centering
\caption{\label{t-urb} Interaction of urban-rural at person level with \% urban
  at country level. X-axis  values are chosen based on the variable's
  distribution: 31.2 is 5th percentile, median is 75, and 90 is 92nd percentile.}
\end{figure}

We have used above just the usual World Bank definition, which uses each
country's statistical office definition--for discussion
see especially \citet{ebshoyEgypt24} and also \citet{deb23}.


We have also tries similar cross-level interactions with PCGDP and income
inequality, but didn't find significant interactions. 



Last we check alternative specification if we are not overcontrolling
\citep{bartram24}--in the body of the paper we used a set of control variables
from the literature, however, we think that we may drop gender, age and age2--as
they may not be predictive of urbanicity. First, we repeat graph from the body
for ease of comparison in figure \ref{dd1}. And below it figure without the
above mentioned controls \ref{dd2}. Results are substantively the same.  And in
\ref{dd3} adding education, eg more educated persons may move to a city--again,
 very similar results. And in \ref{dd4} fewer countries but still most
 insignificant. 


\begin{figure}[H]
 \includegraphics[height=5in]{fig-regCC1mod2}\centering
\caption{\label{dd1}} 
\end{figure}
\begin{figure}[H]
 \includegraphics[height=5in]{fig-regCC1mod2w}\centering
\caption{\label{dd2} Dropping age age2 and gender} 
\end{figure}
\begin{figure}[H]
 \includegraphics[height=5in]{fig-regCC1mod2ww}\centering
\caption{\label{dd3} Adding education--again,
very similar results} 
\end{figure}
\begin{figure}[H]
 \includegraphics[height=5in]{fig-regCC1mod2ww20}\centering
\caption{\label{dd4} only $<$2020--fewer countries but still most insignificant.} 
\end{figure}


\subsubsection{Interactions with financial satisfaction, class, and income}

\begin{figure}[H]
 \includegraphics[height=3in]{m-app.pdf}\centering
\caption{\label{m-app} Multilevel country random intercept regressions of
  freedom on rural-urban dummies interacted with: a) financial satisfaction, b)
  class, c) income. For underlying regressions see appendix. Controls are
  the same as in the previous section.}
\end{figure}


\begin{spacing}{.9} \begin{table}[H]\centering\begin{scriptsize}
\begin{tabular}{p{2.8in}p{.55in}p{.55in}p{.55in}p{.55in}p{.55in}p{.55in}p{.55in}p{.55in}p{.55in}}\hline
\input{regABC.tex}\hline + 0.10 * 0.05 ** 0.01 *** 0.001; robust std err
\end{tabular}\end{scriptsize}\caption{\label{regABC}OLS regressions of freedom/agency}
\end{table}\end{spacing}

We could have explored more fully here using usual sequential elaboration of the
models, but for now just sticking with the final full model as there are already
plenty of robustness/sensitivity checks and one could easily get lost in this
sea of extra results. 

\begin{figure}[H]
 \includegraphics[height=5in]{m-body2}\centering
\caption{\label{m-body2} Robustness/sensitivity check: on the sample after drooping countries where top
  gt500k cat missing. Results are mostly similar to those reported in the body
  of the paper.} 
\end{figure}

Last we check alternative specification if we are not overcontrolling
\citep{bartram24}--in the body of the paper we used a set of control variables
from the literature, however, we think that we may drop gender, age and age2--as
they may not be predictive of urbanicity. Results are in figure \ref{m-body2q} and are
substantively the same, and so when adding education in figure
\ref{m-body2qv}. The very final check is dropping pandemic years 2020 and
later--pandemic might have made city life less free, but our results remain
similar, if anything stronger in figure \ref{m-body2qv20}  

\begin{figure}[H]
 \includegraphics[height=5in]{m-body2q}\centering
\caption{\label{m-body2q} Dropping age age2 and gender.} 
\end{figure}
\begin{figure}[H]
 \includegraphics[height=5in]{m-body2qv}\centering
\caption{\label{m-body2qv} Adding education.} 
\end{figure}
\begin{figure}[H]
 \includegraphics[height=5in]{m-body2qv20}\centering
\caption{\label{m-body2qv20} Only $<$2020.} 
\end{figure}




{\bf Initial interactions and alternative specifications}


\begin{scriptsize}
\begin{spacing}{.19} 
\begin{verbatim}
. mixed free c.town##c.satFin  trust famUns criVic  age age2 male mar i.emp he
> alth kids rel_imp  || cc:   ,mle

Performing EM optimization ...

Performing gradient-based optimization: 
Iteration 0:  Log likelihood =  -184956.1  
Iteration 1:  Log likelihood =  -184956.1  

Computing standard errors ...

Mixed-effects ML regression                        Number of obs    =   86,280
Group variable: cc                                 Number of groups =       64
                                                   Obs per group:
                                                                min =      747
                                                                avg =  1,348.1
                                                                max =    4,018
                                                   Wald chi2(20)    = 11068.66
Log likelihood =  -184956.1                        Prob > chi2      =   0.0000

------------------------------------------------------------------------------
        free | Coefficient  Std. err.      z    P>|z|     [95% conf. interval]
-------------+----------------------------------------------------------------
        town |      -0.03       0.01    -3.99    0.00        -0.05       -0.02
      satFin |       0.22       0.01    31.95    0.00         0.20        0.23
             |
      c.town#|
    c.satFin |       0.01       0.00     5.90    0.00         0.00        0.01
             |
       trust |       0.02       0.02     1.34    0.18        -0.01        0.06
      famUns |      -0.11       0.01   -12.26    0.00        -0.13       -0.09
      criVic |      -0.11       0.03    -4.22    0.00        -0.16       -0.06
         age |       0.02       0.00     5.99    0.00         0.01        0.02
        age2 |      -0.00       0.00    -5.37    0.00        -0.00       -0.00
        male |       0.01       0.02     0.67    0.50        -0.02        0.04
         mar |       0.00       0.02     0.03    0.98        -0.03        0.03
             |
         emp |
  Part time  |      -0.08       0.03    -2.82    0.00        -0.13       -0.02
Self empl..  |      -0.01       0.02    -0.56    0.57        -0.06        0.03
    Retired  |      -0.04       0.03    -1.32    0.19        -0.10        0.02
  Housewife  |      -0.20       0.03    -8.00    0.00        -0.25       -0.15
   Students  |      -0.14       0.04    -3.79    0.00        -0.21       -0.07
 Unemployed  |      -0.10       0.03    -3.29    0.00        -0.15       -0.04
      Other  |      -0.22       0.07    -3.28    0.00        -0.36       -0.09
             |
      health |       0.33       0.01    37.13    0.00         0.31        0.35
        kids |       0.02       0.01     2.67    0.01         0.00        0.03
     rel_imp |       0.04       0.01     4.11    0.00         0.02        0.05
       _cons |       4.26       0.11    38.13    0.00         4.04        4.48
------------------------------------------------------------------------------

------------------------------------------------------------------------------
  Random-effects parameters  |   Estimate   Std. err.     [95% conf. interval]
-----------------------------+------------------------------------------------
cc: Identity                 |
                  var(_cons) |       0.28       0.05          0.19        0.39
-----------------------------+------------------------------------------------
               var(Residual) |       4.25       0.02          4.21        4.29
------------------------------------------------------------------------------
LR test vs. linear model: chibar2(01) = 4378.78       Prob >= chibar2 = 0.0000
\end{verbatim}
\end{spacing}{.19}
\end{scriptsize}
 \begin{figure}[H]
  \includegraphics[height=3in]{ini1.pdf}\centering
 \caption{\label{ini1}ini1}
 \end{figure}
few things to note in the above--we have trust, famUns and criVic--which are
rather predicted by urbanicity, not predictors, hence we may have overcontrol
bias. Also here both town and satFin are forced to be linear, a tenuous
assumption. 

\begin{scriptsize}
\begin{spacing}{.19} 
\begin{verbatim}
 mixed free i.t4##i.sf3  trust famUns criVic  age age2 male mar i.emp health 
> kids rel_imp  || cc:   ,mle

Performing EM optimization ...

Performing gradient-based optimization: 
Iteration 0:  Log likelihood = -185241.39  
Iteration 1:  Log likelihood = -185241.39  (backed up)

Computing standard errors ...

Mixed-effects ML regression                        Number of obs    =   86,280
Group variable: cc                                 Number of groups =       64
                                                   Obs per group:
                                                                min =      747
                                                                avg =  1,348.1
                                                                max =    4,018
                                                   Wald chi2(28)    = 10427.08
Log likelihood = -185241.39                        Prob > chi2      =   0.0000

------------------------------------------------------------------------------
        free | Coefficient  Std. err.      z    P>|z|     [95% conf. interval]
-------------+----------------------------------------------------------------
          t4 |
    10-100k  |      -0.14       0.04    -3.39    0.00        -0.21       -0.06
   100-500k  |      -0.04       0.05    -0.78    0.43        -0.13        0.06
     gt500k  |      -0.03       0.04    -0.74    0.46        -0.12        0.05
             |
         sf3 |
          2  |       0.65       0.03    18.99    0.00         0.59        0.72
          3  |       1.50       0.04    40.61    0.00         1.42        1.57
             |
      t4#sf3 |
  10-100k#2  |       0.21       0.05     4.38    0.00         0.12        0.31
  10-100k#3  |       0.12       0.05     2.38    0.02         0.02        0.22
 100-500k#2  |       0.09       0.06     1.62    0.11        -0.02        0.20
 100-500k#3  |       0.07       0.06     1.17    0.24        -0.05        0.19
   gt500k#2  |       0.12       0.05     2.43    0.01         0.02        0.22
   gt500k#3  |       0.12       0.05     2.26    0.02         0.02        0.23
             |
       trust |       0.03       0.02     1.80    0.07        -0.00        0.07
      famUns |      -0.12       0.01   -13.44    0.00        -0.14       -0.10
      criVic |      -0.12       0.03    -4.75    0.00        -0.17       -0.07
         age |       0.02       0.00     5.32    0.00         0.01        0.02
        age2 |      -0.00       0.00    -4.66    0.00        -0.00       -0.00
        male |       0.01       0.02     0.70    0.48        -0.02        0.04
         mar |       0.01       0.02     0.34    0.73        -0.03        0.04
             |
         emp |
  Part time  |      -0.08       0.03    -3.06    0.00        -0.14       -0.03
Self empl..  |      -0.02       0.02    -0.70    0.48        -0.06        0.03
    Retired  |      -0.05       0.03    -1.66    0.10        -0.11        0.01
  Housewife  |      -0.22       0.03    -8.41    0.00        -0.27       -0.17
   Students  |      -0.14       0.04    -3.88    0.00        -0.21       -0.07
 Unemployed  |      -0.13       0.03    -4.41    0.00        -0.19       -0.07
      Other  |      -0.25       0.07    -3.65    0.00        -0.39       -0.12
             |
      health |       0.35       0.01    39.72    0.00         0.34        0.37
        kids |       0.01       0.01     2.26    0.02         0.00        0.03
     rel_imp |       0.04       0.01     4.41    0.00         0.02        0.06
       _cons |       4.82       0.11    45.42    0.00         4.61        5.03
------------------------------------------------------------------------------

------------------------------------------------------------------------------
  Random-effects parameters  |   Estimate   Std. err.     [95% conf. interval]
-----------------------------+------------------------------------------------
cc: Identity                 |
                  var(_cons) |       0.27       0.05          0.19        0.39
-----------------------------+------------------------------------------------
               var(Residual) |       4.28       0.02          4.23        4.32
------------------------------------------------------------------------------
LR test vs. linear model: chibar2(01) = 4250.93       Prob >= chibar2 = 0.0000
\end{verbatim}
\end{spacing}{.19}
\end{scriptsize}
 \begin{figure}[H]
  \includegraphics[height=3in]{ini2.pdf}\centering
 \caption{\label{ini2}ini2}
 \end{figure}
In the above a different take--financial satisfaction tertiles--for the 1st
tertile the biggest drop is actually in the middle category at 10-100k.
 
\begin{scriptsize}
\begin{spacing}{.19} 
\begin{verbatim}
 mixed free c.town##c.inc  trust famUns criVic  age age2 male mar i.emp healt
> h kids rel_imp  || cc:   ,mle

Performing EM optimization ...

Performing gradient-based optimization: 
Iteration 0:  Log likelihood = -184725.92  
Iteration 1:  Log likelihood = -184725.92  (backed up)

Computing standard errors ...

Mixed-effects ML regression                         Number of obs    =  84,979
Group variable: cc                                  Number of groups =      64
                                                    Obs per group:
                                                                 min =     749
                                                                 avg = 1,327.8
                                                                 max =   4,018
                                                    Wald chi2(20)    = 5460.49
Log likelihood = -184725.92                         Prob > chi2      =  0.0000

------------------------------------------------------------------------------
        free | Coefficient  Std. err.      z    P>|z|     [95% conf. interval]
-------------+----------------------------------------------------------------
        town |      -0.00       0.01    -0.43    0.67        -0.02        0.01
         inc |       0.11       0.01    13.01    0.00         0.09        0.12
             |
c.town#c.inc |       0.00       0.00     1.21    0.23        -0.00        0.00
             |
       trust |       0.07       0.02     3.91    0.00         0.04        0.11
      famUns |      -0.17       0.01   -18.82    0.00        -0.19       -0.16
      criVic |      -0.16       0.03    -6.05    0.00        -0.21       -0.11
         age |       0.01       0.00     2.60    0.01         0.00        0.01
        age2 |      -0.00       0.00    -1.15    0.25        -0.00        0.00
        male |      -0.00       0.02    -0.17    0.87        -0.03        0.03
         mar |       0.04       0.02     1.97    0.05         0.00        0.07
             |
         emp |
  Part time  |      -0.08       0.03    -2.83    0.00        -0.14       -0.02
Self empl..  |      -0.01       0.02    -0.46    0.65        -0.06        0.04
    Retired  |       0.01       0.03     0.33    0.74        -0.05        0.07
  Housewife  |      -0.18       0.03    -6.65    0.00        -0.23       -0.12
   Students  |      -0.11       0.04    -2.99    0.00        -0.18       -0.04
 Unemployed  |      -0.15       0.03    -5.04    0.00        -0.21       -0.09
      Other  |      -0.25       0.07    -3.54    0.00        -0.39       -0.11
             |
      health |       0.46       0.01    51.21    0.00         0.45        0.48
        kids |       0.01       0.01     1.93    0.05        -0.00        0.03
     rel_imp |       0.05       0.01     5.88    0.00         0.04        0.07
       _cons |       4.75       0.12    39.88    0.00         4.51        4.98
------------------------------------------------------------------------------

------------------------------------------------------------------------------
  Random-effects parameters  |   Estimate   Std. err.     [95% conf. interval]
-----------------------------+------------------------------------------------
cc: Identity                 |
                  var(_cons) |       0.35       0.06          0.25        0.50
-----------------------------+------------------------------------------------
               var(Residual) |       4.51       0.02          4.47        4.55
------------------------------------------------------------------------------
LR test vs. linear model: chibar2(01) = 5398.66       Prob >= chibar2 = 0.0000
\end{verbatim}
\end{spacing}{.19}
\end{scriptsize}
 \begin{figure}[H]
  \includegraphics[height=3in]{ini3.pdf}\centering
 \caption{\label{ini3}ini3}
 \end{figure}
 Here what is easy to see is that along the income spectrum
 the largest places $>500k$ change most; and
smallest places $<2k$ change the least
 

 
\begin{scriptsize}
\begin{spacing}{.19} 
\begin{verbatim}
 mixed free c.town##c.class  trust famUns criVic  age age2 male mar i.emp hea
> lth kids rel_imp  || cc:   ,mle

Performing EM optimization ...

Performing gradient-based optimization: 
Iteration 0:  Log likelihood = -184877.86  
Iteration 1:  Log likelihood = -184877.86  

Computing standard errors ...

Mixed-effects ML regression                         Number of obs    =  84,931
Group variable: cc                                  Number of groups =      64
                                                    Obs per group:
                                                                 min =     686
                                                                 avg = 1,327.0
                                                                 max =   4,018
                                                    Wald chi2(20)    = 4850.01
Log likelihood = -184877.86                         Prob > chi2      =  0.0000

------------------------------------------------------------------------------
        free | Coefficient  Std. err.      z    P>|z|     [95% conf. interval]
-------------+----------------------------------------------------------------
        town |      -0.00       0.01    -0.52    0.60        -0.02        0.01
       class |       0.14       0.02     7.62    0.00         0.10        0.18
             |
      c.town#|
     c.class |       0.01       0.00     1.56    0.12        -0.00        0.01
             |
       trust |       0.09       0.02     4.79    0.00         0.05        0.13
      famUns |      -0.18       0.01   -19.22    0.00        -0.20       -0.16
      criVic |      -0.17       0.03    -6.34    0.00        -0.22       -0.11
         age |       0.01       0.00     2.45    0.01         0.00        0.01
        age2 |      -0.00       0.00    -1.24    0.22        -0.00        0.00
        male |       0.01       0.02     0.41    0.68        -0.02        0.04
         mar |       0.05       0.02     2.96    0.00         0.02        0.09
             |
         emp |
  Part time  |      -0.10       0.03    -3.45    0.00        -0.16       -0.04
Self empl..  |      -0.03       0.02    -1.22    0.22        -0.08        0.02
    Retired  |      -0.02       0.03    -0.64    0.52        -0.08        0.04
  Housewife  |      -0.21       0.03    -7.84    0.00        -0.26       -0.16
   Students  |      -0.13       0.04    -3.46    0.00        -0.20       -0.06
 Unemployed  |      -0.21       0.03    -6.93    0.00        -0.27       -0.15
      Other  |      -0.30       0.07    -4.13    0.00        -0.44       -0.16
             |
      health |       0.48       0.01    52.29    0.00         0.46        0.49
        kids |       0.01       0.01     1.81    0.07        -0.00        0.02
     rel_imp |       0.05       0.01     5.88    0.00         0.04        0.07
       _cons |       4.86       0.12    39.41    0.00         4.61        5.10
------------------------------------------------------------------------------

------------------------------------------------------------------------------
  Random-effects parameters  |   Estimate   Std. err.     [95% conf. interval]
-----------------------------+------------------------------------------------
cc: Identity                 |
                  var(_cons) |       0.36       0.06          0.25        0.51
-----------------------------+------------------------------------------------
               var(Residual) |       4.54       0.02          4.49        4.58
------------------------------------------------------------------------------
LR test vs. linear model: chibar2(01) = 5460.71       Prob >= chibar2 = 0.0000
\end{verbatim}
\end{spacing}{.19}
\end{scriptsize}
 \begin{figure}[H]
  \includegraphics[height=3in]{ini4.pdf}\centering
 \caption{\label{ini4}ini4}
 \end{figure}
The above is a version of the main graph from the body of the paper, again here
we are over-controlling with trust and the 2 crime variables, but it is not that
different after all, similar patters, just the main difference being that here
the lower class is flat and in the body downward sloping and here the top class
is quite steep and while in paper it is upward sloping it's little less so in
the body of the paper

 \begin{figure}[H]
  \includegraphics[height=3in]{ini5.pdf}\centering
 \caption{\label{ini5}ini5}
 \end{figure}
ini5 is from the same model, nicely shows that the higher the class the more
difference between the smallest and the largest places


\begin{scriptsize}
\begin{spacing}{.19} 
\begin{verbatim}
 mixed free c.town##i.class  trust famUns criVic  age age2 male mar i.emp hea
> lth kids rel_imp  || cc:   ,mle

Performing EM optimization ...

Performing gradient-based optimization: 
Iteration 0:  Log likelihood = -184853.93  
Iteration 1:  Log likelihood = -184853.93  (backed up)

Computing standard errors ...

Mixed-effects ML regression                         Number of obs    =  84,931
Group variable: cc                                  Number of groups =      64
                                                    Obs per group:
                                                                 min =     686
                                                                 avg = 1,327.0
                                                                 max =   4,018
                                                    Wald chi2(24)    = 4900.64
Log likelihood = -184853.93                         Prob > chi2      =  0.0000

------------------------------------------------------------------------------
        free | Coefficient  Std. err.      z    P>|z|     [95% conf. interval]
-------------+----------------------------------------------------------------
        town |      -0.00       0.01    -0.54    0.59        -0.02        0.01
             |
       class |
Working c..  |       0.22       0.06     3.86    0.00         0.11        0.34
Lower mid..  |       0.28       0.06     5.10    0.00         0.17        0.39
Upp mid/u..  |       0.48       0.06     7.75    0.00         0.36        0.60
             |
class#c.town |
Working c..  |       0.01       0.01     1.34    0.18        -0.01        0.03
Lower mid..  |       0.01       0.01     1.28    0.20        -0.01        0.03
Upp mid/u..  |       0.02       0.01     1.95    0.05        -0.00        0.04
             |
       trust |       0.09       0.02     4.82    0.00         0.05        0.13
      famUns |      -0.18       0.01   -19.18    0.00        -0.20       -0.16
      criVic |      -0.16       0.03    -6.30    0.00        -0.22       -0.11
         age |       0.01       0.00     2.50    0.01         0.00        0.01
        age2 |      -0.00       0.00    -1.29    0.20        -0.00        0.00
        male |       0.01       0.02     0.35    0.72        -0.03        0.04
         mar |       0.05       0.02     2.87    0.00         0.02        0.09
             |
         emp |
  Part time  |      -0.10       0.03    -3.39    0.00        -0.15       -0.04
Self empl..  |      -0.03       0.02    -1.10    0.27        -0.07        0.02
    Retired  |      -0.02       0.03    -0.52    0.60        -0.08        0.05
  Housewife  |      -0.20       0.03    -7.60    0.00        -0.25       -0.15
   Students  |      -0.13       0.04    -3.34    0.00        -0.20       -0.05
 Unemployed  |      -0.20       0.03    -6.63    0.00        -0.26       -0.14
      Other  |      -0.29       0.07    -4.01    0.00        -0.43       -0.15
             |
      health |       0.47       0.01    52.12    0.00         0.46        0.49
        kids |       0.01       0.01     1.85    0.06        -0.00        0.03
     rel_imp |       0.05       0.01     5.85    0.00         0.04        0.07
       _cons |       4.96       0.12    40.78    0.00         4.72        5.20
------------------------------------------------------------------------------
\end{verbatim}
\end{spacing}{.19}
\end{scriptsize}
 \begin{figure}[H]
  \includegraphics[height=3in]{ini6.pdf}\centering
 \caption{\label{ini6}ini6}
 \end{figure}
 And then ini6: couple interesting things: first things flip from lower to
 working class, then 
 working and lower middle class quite similar, and the differences get larger
 for upper middle and upper classes. 
 
 
% \begin{scriptsize}
% \begin{spacing}{.19} 
% \begin{verbatim}

% \end{verbatim}
% \end{spacing}{.19}
% \end{scriptsize}
%  \begin{figure}[H]
%   \includegraphics[height=3in]{??.pdf}\centering
%  \caption{\label{??}??}
%  \end{figure}




 
\subsubsection{By region and country income group: Interactions with financial
  satisfaction} %class income

Here we repeat urbanicity interactions with financial satisfaction from the body
of the paper but by region and income group; the first panel in each graph is
repeated from the body of the paper.  

In figure \ref{zzz} we see that the pattern is much stronger for East
Asia/Pacific, and opposite especially in Middle East/North Africa, but also in
North America. 

 \begin{figure}[H]
  \includegraphics[height=3in]{zzz.pdf}\centering
 \caption{\label{zzz} Urbanicity interactions with financial satisfaction by region.}
 \end{figure}

 In figure \ref{xxx} high income has a similar pattern,
 and lower middle income panel has a strong pattern, and quite the
opposite pattern for upper middle income.  

 \begin{figure}[H]
  \includegraphics[height=3in]{xxx.pdf}\centering
  \caption{\label{xxx} Urbanicity interactions with financial satisfaction by
    income group. Note just 2 countries for low income group: ETH TJK.}
 \end{figure}


 



\section{Auxiliary results and sidenote} 

\subsection{Some graphs}
%/home/aok/papers/leonieAgency/tex/ 
fig \ref{free_ine} and \ref{free_pov} 

weird to think that more  ineq and pov  would produce more freedom/autonomy;
rather i guess lax capitalist friendly regulation would produce more 
ineq and pov, but should also produce less autonomy/freedom on avg, but it
doesn't, there is positive relationship\\

or maybe the other way round: more freedom/autonomy more ineq and pov




 \begin{figure}[H]
  \includegraphics[height=3in]{/home/aok/papers/leonieAgency/tex/free_ine.pdf}\centering
 \caption{\label{free_ine}more ineq (inc share held by top 10perc, more autonomy--conservatives
   would like it!! and pretty strong corr .4}
 \end{figure}

  \begin{figure}[H]
  \includegraphics[height=3in]{/home/aok/papers/leonieAgency/tex/free_pov.pdf}\centering
 \caption{\label{free_pov}likewise with poverty, but here corr just .2}
 \end{figure}



\subsection{Literature}

\citet{newmanMisc20oct27} makes some useful points that may be worth citing in full:
 \begin{quote}
   Social theory typically conceives of urban growth as associated with
   political freedom. [...]  But not even in the most likely of cases did the
   city makes urban peasants free. Focusing on Mexico City, I argue that
   because rural-to-urban migrants became urban denizens by establishing
   squatter settlements, they entered into dependent relations with political
   elites, which reinforced a system of political dependency, or clientelist
   politics. [...] Mexico is, perhaps, the most iconic case. A large cross
   section of the peasantry rose in arms against landowning oligarchs in the
   insurrectionary phase of the Mexican Revolution (1910-1917), ultimately
   precipitating the wholesale reconstruction of the political order.
   Thereafter, peasants continued to mobilize for radical change, constituting a
   major obstacle to postrevolutionary governing elites who sought to
   consolidate control. [...] When a large and growing portion of rural-to-urban
   migrants live in squatter settlements, and therefore have little recourse to
   legal protection in residential matters, they almost inevitably grow
   dependent on political elites--in an arrangement known as ``clientelism.'' Urban
   concentration gives rise to clientelism.  
 \end{quote}



 
% is organized byu section, not subsection
% !!!
% have most of the stuff outputted to online appendix:)--start with that and then
% select stuff to paper--have brief narrative describng patterns in online app too
% !!!

% \section*{Variables' definitions, coding, and distributions}
% \label{app_var_des}


% %\input{/tmp/a.tex} %aok_var_des

% % \begin{spacing}{.9}
% %   \begin{table}[H]\centering \caption{Summary statistics.} \label{sumSta} \begin{scriptsize} \begin{tabular}{p{1.8in}p{.5in}p{.5in}p{.5in}p{.5in}p{.5in}p{.5in}p{.5in}p{.5in}p{.5in}p{.5
% %             in}p{.5in}p{.5 in}}\hline
% %         \input{/tmp/aha2.tex}
% %          \end{tabular}\end{scriptsize}\end{table}
% % \end{spacing}

% % \begin{spacing}{.9}
% %   \begin{table}[H]\centering \caption{Correlation matrix.} \label{sumSta} \begin{scriptsize} \begin{tabular}{@{}
% %           p{1.2in} rrrrrrrrrrrrr @{}}\hline
% %         \input{/tmp/ahb2.tex}\hline
% %          \end{tabular}\end{scriptsize}\end{table}
% % \end{spacing}



% Table XXX shows variable distributions. If a variable has more than
% 10 categories it is classified into bins...

% %\input .... %TODO !!!! have input here histograms

% \section*{Additional Descriptive Statistics}
% \label{app_des_sta}

% %make sure i have [H] or h! ???
% % \begin{table}[H]
% % \caption{}
% % \centering
% % \label{}
% % \begin{scriptsize}
% % \input{../out/reg_c.tex}
% % \end{scriptsize}
% % \end{table}

%\newpage
%\theendnotes


\begin{scriptsize}
\begin{spacing}{.9} 
\begin{verbatim}
. tabstat ls, by(lwd7) stat(mean sd) format(%9.1f)

Summary for variables: ls
Group variable: lwd7 (RECODE of lwd (do what you like))

  lwd7 |      Mean        SD
-------+--------------------
   0-4 |       5.6       2.7
     2 |       7.6       2.3
     3 |       7.6       1.8
     4 |       8.0       1.5
     5 |       7.8       1.9
     6 |       8.2       1.5
     7 |       8.9       1.5
-------+--------------------
 Total |       8.3       1.9
----------------------------
\end{verbatim}
\end{spacing}{.9}
\end{scriptsize}


\bibliography{/home/aok/papers/root/tex/ebib.bib,cityAgency.bib,/home/aok/papers/root/old/2021/swbCityWorld/tex/swbCityWorld.bib}

\end{spacing}
\end{document}
